\chapter{协同搬运与多机力控}
\label{ch:mr-coopforce}

% ════════════════════════════════════════════════════════════════
%  全吸收章：重渲染 Tutorial 笔记
%   - 04_移动机器人规控/50_多机器人协作/60_协同搬运与力控.md
%     （闭链/抓取矩阵 / 内力外力分解 / leader-follower / 去中心化 / 物体阻抗 /
%      负载分配 QP / 无源性·能量罐·波变量 / 绳索·被动臂·零通信学习前沿）
%  记号归一：R∈SO(3)、T∈SE(3)（preamble 宏 \SO \SE \so \se）；力学/接触/优化自然记号保留。
%  源由 AI 生成，作者归属/年份/数值/前沿论文编号存疑处就地 \pz/\rebuilt，并入 claims 台账。
%  教学层遵循 v5.0 理论教学规范（动机→反面→历史→理论→陷阱→练习 + 误解汇总/小结/速查/故障排查）。
% ════════════════════════════════════════════════════════════════

% —— 本章局部数学算子（章文件 \input 入正文，故不用 \DeclareMathOperator，
%    改用 \providecommand + \operatorname 兜底，避免与 preamble 既有定义冲突）——
\providecommand{\rank}{\operatorname{rank}}
\providecommand{\range}{\operatorname{range}}
\providecommand{\rowsp}{\operatorname{row}}
\providecommand{\diagop}{\operatorname{diag}}
\providecommand{\blkdiag}{\operatorname{blkdiag}}
\providecommand{\trace}{\operatorname{tr}}
% 反对称（叉积）矩阵：本章接触点位置向量 r 的 \widehat{r}（与 \cref{ch:lie} 的 (·)^\wedge 同义）
\providecommand{\skewm}[1]{\widehat{#1}}

\rebuilt{本章重渲染自 Tutorial 笔记《协同搬运与力控》，其源稿由 AI 撰写，论文归属、年份、卷期与前沿工作（如 Walker 非挤压伪逆、Force-ANTS、PACC、decPLM）的编号、数值可能有误。本章忠实重述、不擅自“修正”，逐处 \texttt{\textbackslash cite} 标源；存疑断言已登记 claims 台账，待联网核对后二次校验。}

\begin{practice}[前置自测]
开始之前，先试答下列问题。若已能流畅作答，可只速读本章表格与速查卡；若卡住（答不出 $\ge 2$ 题），第 1、2 题回操作空间动力学与力控导论、第 3 题回数值线性代数（伪逆/零空间）、第 4 题回多机系统全景与共识、第 5 题回无源性基础，再来读本章——本章每一步推导都建立在这些之上。
\begin{enumerate}
  \item 单臂操作空间动力学 $\boldsymbol{\Lambda}(\mathbf{x})\ddot{\mathbf{x}}+\boldsymbol{\mu}+\mathbf{p}=\mathbf{F}$ 中，操作空间惯量 $\boldsymbol{\Lambda}=(\mathbf{J}\mathbf{M}^{-1}\mathbf{J}^\top)^{-1}$ 的物理含义是什么？雅可比转置力映射 $\boldsymbol{\tau}=\mathbf{J}^\top\mathbf{F}$ 怎么由虚功原理得到？
  \item Hogan 阻抗律 $\mathbf{F}=\mathbf{M}_d\ddot{\tilde{\mathbf{x}}}+\mathbf{D}_d\dot{\tilde{\mathbf{x}}}+\mathbf{K}_d\tilde{\mathbf{x}}$ 三个矩阵各整形什么物理行为？阻抗与导纳的因果方向有何区别？为什么刚性环境下导纳易不稳？
  \item 给定“宽矩阵” $\mathbf{G}\in\mathbb{R}^{6\times 3N}$（$3N>6$），$\mathbf{G}\mathbf{f}=\mathbf{w}$ 的解集是什么结构？Moore--Penrose 伪逆 $\mathbf{G}^{+}=\mathbf{G}^\top(\mathbf{G}\mathbf{G}^\top)^{-1}$ 给出的特解有何最优性？零空间投影 $\mathbf{I}-\mathbf{G}^{+}\mathbf{G}$ 把任意向量投到哪里？
  \item 什么叫去中心化（decentralized）控制？为什么“每个机器人只用本地测量、不通信”在多机搬运里反而是优点？隐式通信通过什么物理量传递信息？
  \item 一个系统“无源（passive）”的能量定义是什么？为什么无源蕴含稳定？两个无源系统反馈互联后还无源吗？
\end{enumerate}
\noindent\emph{答案线索}：1、2 回单臂力控章（操作空间动力学、阻抗/导纳二分法）；3 回 \nameref{ch:notation} 与数值线性代数（伪逆与零空间）；4 回多机系统全景与 \cref{ch:mr-consensus}（共识）；5 回无源性基础——这五点正是本章 \cref{sec:cm-grasp,sec:cm-internal,sec:cm-impedance,sec:cm-decentral,sec:cm-passivity} 的地基。
\end{practice}

\section*{本章目标}
学完本章，你应当能够：
\begin{enumerate}
  \item \textbf{从闭链约束推导抓取矩阵}：把 $N$ 个机器人刚性抓持同一物体建模为闭合运动链，写出抓取矩阵 $\mathbf{G}$ 的解析结构（每列 $[\mathbf{I}_3;\ \skewm{\mathbf{r}}_i]$ 的来历），解释力映射 $\mathbf{w}=\mathbf{G}\mathbf{f}$ 与速度对偶 $\dot{\mathbf{x}}_i=\mathbf{G}_i^\top\mathbf{v}_o$，并说明零空间维度为何恒为 $3N-6$；
  \item \textbf{严格分解内力与外力}：证明接触力分解 $\mathbf{f}=\mathbf{G}^{+}\mathbf{w}+(\mathbf{I}-\mathbf{G}^{+}\mathbf{G})\boldsymbol{\eta}$ 两项的正交性与物理意义，理解为什么 Moore--Penrose 解并非“非挤压（nonsqueezing）”解，掌握 Walker 非挤压伪逆与内力基构造；
  \item \textbf{对比 leader--follower 与去中心化两条路线}：讲清主从架构（leader 用阻抗领航、follower 估计意图）与全去中心化架构（同构 agent、物体隐式协调）各自的假设、信息流、可扩展性与失效模式；
  \item \textbf{设计物体阻抗控制器}：推导 object impedance（把期望阻抗加在物体质心而非各末端）的控制律，理解它与“各末端独立阻抗”的区别，掌握把物体级期望力旋量经 $\mathbf{G}^{+}$ 与内力调节分配回各机器人的完整管线；
  \item \textbf{做负载分配优化}：把外力旋量分配形式化为带摩擦锥/单边/力限约束的二次规划（QP），理解等分、加权伪逆、QP 三种策略的取舍；
  \item \textbf{用无源性判稳并保稳}：以 Lyapunov/无源性框架分析多机搬运闭环，理解能量罐（energy tank）、波变量（wave variable）如何在变阻抗与通信延迟下保住无源性；
  \item \textbf{打通知识网络}：明确本章抓取矩阵理论与双臂协调力控、操作空间动力学、单臂无源性是同一套数学在不同 agent 数下的展开，能在多足绳索搬运、被动臂搬运、零通信学习搬运等前沿之间建立联系。
\end{enumerate}

\subsection*{本章知识导航}
\textbf{一句话定位}：当多台机器人同时抓住同一刚体一起搬运时，它们被这个物体\emph{焊死成一条闭合运动链}——任何一台多用一分力，这分力都不会让物体动得更快，而是变成挤压物体的“内力”。本章从这条闭链约束出发，讲透\emph{抓取矩阵如何把接触力映射成物体合力旋量、内力/外力如何沿零空间分解、leader--follower 与去中心化两条协调路线的本质差异、阻抗如何在多机之间协调、以及为什么无源性是判断这套耦合系统稳不稳的终极标尺}。

\begin{center}
\small
\begin{tikzpicture}[
  node distance=5mm and 7mm, >=Stealth,
  blk/.style={draw, rounded corners, align=center, font=\small, inner sep=3pt, fill=main!6, draw=main!60!black},
  grp/.style={draw, rounded corners, align=center, font=\small, inner sep=3pt, fill=second!6, draw=second!60!black},
  fin/.style={draw, rounded corners, align=center, font=\small, inner sep=3pt, fill=third!8, draw=third!60!black},
  ln/.style={->, thick, gray!60!black}]
\node[blk] (grasp) {闭链 $+$ 抓取矩阵\\ $\mathbf{w}=\mathbf{G}\mathbf{f}$、$\dot{\mathbf{x}}_i=\mathbf{G}_i^\top\mathbf{v}_o$};
\node[blk, right=of grasp] (intf) {内力 / 外力分解\\ $\ker\mathbf{G}$ 维度 $3N-6$};
\node[grp, below=9mm of grasp] (lead) {leader--follower\\ 阻抗领航 $+$ 意图估计};
\node[grp, right=of lead] (dec) {去中心化\\ 物体即共享黑板};
\node[blk, below=9mm of lead] (imp) {物体阻抗\\ 质心柔顺 $+$ 内力解耦};
\node[blk, right=of imp] (qp) {负载分配 QP\\ 摩擦锥 / 单边 / 力限};
\node[fin, below=9mm of imp] (pass) {无源性与稳定性\\ 能量罐、波变量、互联封闭};
\node[fin, right=of pass] (front) {前沿与综合\\ 绳索 / 被动臂 / 零通信学习};
\draw[ln] (grasp)--(intf);
\draw[ln] (grasp.south)--(lead.north);
\draw[ln] (lead)--(dec);
\draw[ln] (intf.south) -- ++(0,-3mm) -| (dec.north);
\draw[ln] (lead.south)--(imp.north);
\draw[ln] (imp)--(qp);
\draw[ln] (imp.south)--(pass.north);
\draw[ln] (pass)--(front);
\end{tikzpicture}
\end{center}

\noindent\textbf{推荐路径}：\cref{sec:cm-grasp}（闭链/抓取矩阵）与 \cref{sec:cm-internal}（内力/外力分解）是全章地基，任何读者必读、需跟着推一遍；\cref{sec:cm-leader,sec:cm-decentral}（两条协调路线）回答“谁来决策”；\cref{sec:cm-impedance}（物体阻抗）把单臂阻抗推广到物体；\cref{sec:cm-qp}（负载分配 QP）把伪逆升级为带约束优化；\cref{sec:cm-passivity}（无源性）是唯一的研究级硬核小节，把零散控制律焊成稳定系统，值得花最多时间；\cref{sec:cm-frontier}（前沿与综合）用经典框架解读最新系统。\textbf{两条阅读线}——理论线 \cref{sec:cm-grasp} $\to$ \cref{sec:cm-internal} $\to$ \cref{sec:cm-impedance} $\to$ \cref{sec:cm-passivity}（重抓取矩阵、内力分解、物体阻抗、无源性证明）；工程线 \cref{sec:cm-grasp} $\to$ \cref{sec:cm-internal} $\to$ \cref{sec:cm-decentral} $\to$ \cref{sec:cm-qp} $\to$ \cref{sec:cm-frontier}（重力分解、去中心化、QP 力分配、综合）。

\subsection*{前置知识桥接}
本章紧接多机系统的离散决策与连续协调（\cref{ch:mr-consensus} 共识与分布式优化、编队与分布式 MPC），并大量复用单臂力控的操作空间与阻抗工具。这里把最关键的前置点重新激活——你不必翻回去就能跟上：
\begin{itemize}
  \item \textbf{从“信息耦合”到“物理耦合”}。\cref{ch:mr-consensus} 等章的多机器人通过\emph{通信图}间接相连——你给我发消息我才知道你的状态；本章的多机器人通过\emph{被搬刚体}直接相连——我推物体，你在你的接触点立刻感受到。被搬物体既是\emph{约束}（闭链）、又是\emph{隐式通信信道}。这是本章与前几章在耦合机理上的本质不同，也是“零通信搬运”（\cref{sec:cm-decentral,sec:cm-frontier}）成立的根。
  \item \textbf{雅可比 $\leftrightarrow$ 抓取矩阵的同构}。操作空间动力学把单臂动力学从关节空间投到末端：操作空间惯量 $\boldsymbol{\Lambda}=(\mathbf{J}\mathbf{M}^{-1}\mathbf{J}^\top)^{-1}$、雅可比转置力映射 $\boldsymbol{\tau}=\mathbf{J}^\top\mathbf{F}$。本章抓取矩阵 $\mathbf{G}$ 与雅可比 $\mathbf{J}$ 是\emph{同一数学对象在不同层级的化身}：$\mathbf{J}$ 把单臂关节速度映到末端速度，$\mathbf{G}^\top$ 把物体速度映到各接触点速度；“雅可比零空间 $=$ 冗余自运动”在本章原样变成“抓取矩阵零空间 $=$ 内力”。
  \item \textbf{阻抗/导纳与无源性}。单臂力控建立了阻抗（运动$\to$力）/导纳（力$\to$运动）二分法，以及变阻抗的无源性危机与能量罐解法。本章 \cref{sec:cm-impedance} 把单臂阻抗推广到“物体阻抗”，\cref{sec:cm-passivity} 把单臂能量罐推广到多机耦合——核心工具（储能函数、$\dot V\le\mathbf{s}^\top\mathbf{f}$、能量罐放行阀）完全一致，只是 agent 从 1 个变成 $N$ 个，耦合通道从“机器人–环境”变成“机器人–物体–机器人”。
  \item \textbf{伪逆与零空间投影}。欠定线性方程 $\mathbf{G}\mathbf{f}=\mathbf{w}$（$\mathbf{G}$ 行满秩）的通解 $=$ 最小范数特解 $\mathbf{G}^{+}\mathbf{w}$（落在行空间 $\rowsp\mathbf{G}$）$+$ 齐次解 $(\mathbf{I}-\mathbf{G}^{+}\mathbf{G})\boldsymbol{\eta}$（落在零空间 $\ker\mathbf{G}$）。本章这两个对象会以完全相同的形式出现，物理含义从“最小范数解”变成“运动力 $+$ 内力”。
\end{itemize}

\begin{note}
\textbf{本章记号与约定.}\;
旋转矩阵记 $\mathbf{R}\in\SO(3)$、位姿 $\mathbf{T}\in\SE(3)$（preamble 宏 $\SO,\SE,\so,\se$，与全书一致，\nameref{ch:notation}）；力学/接触/优化领域的自然记号保留。
$N$ 为协同搬运的机器人（接触点）数；物体（object）记 $\mathcal{O}$，质心位姿 $\mathbf{x}_o$、广义速度 $\mathbf{v}_o=[\dot{\mathbf{p}}_o^\top,\boldsymbol{\omega}_o^\top]^\top\in\mathbb{R}^6$。
作用在物体质心的合力旋量（wrench）记 $\mathbf{w}_o=[\mathbf{F}_o^\top,\boldsymbol{\tau}_o^\top]^\top\in\mathbb{R}^6$；第 $i$ 接触点的接触力 $\mathbf{f}_i\in\mathbb{R}^3$，堆叠 $\mathbf{f}=[\mathbf{f}_1^\top,\dots,\mathbf{f}_N^\top]^\top\in\mathbb{R}^{3N}$。
$\mathbf{r}_i=\mathbf{x}_i-\mathbf{p}_o$ 为质心到第 $i$ 接触点的位置向量，$\skewm{\mathbf{r}}_i$ 为其反对称（叉积）矩阵，满足 $\skewm{\mathbf{r}}_i\mathbf{v}=\mathbf{r}_i\times\mathbf{v}$（即 \cref{ch:lie} 的 $(\cdot)^\wedge$）。
\textbf{抓取矩阵约定}：全章统一取“$\mathbf{G}$ 把接触力映到物体旋量”，即力映射 $\mathbf{w}_o=\mathbf{G}\mathbf{f}$、速度对偶 $\dot{\mathbf{x}}_i=\mathbf{G}_i^\top\mathbf{v}_o$；与双臂协调力控一致。注意部分文献用 $\mathbf{G}$ 表示该映射的转置（$\mathbf{f}=\mathbf{G}^\top\boldsymbol{\lambda}$ 风格），引用处显式转换。
\textbf{默认接触模型}为点接触带摩擦（每点 $3$ 维力，$\mathbf{G}\in\mathbb{R}^{6\times 3N}$）；涉及双臂刚性抓持时切换到 $6$ 维力旋量模型（$\mathbf{G}\in\mathbb{R}^{6\times 6N}$），就地声明。
\end{note}

\subsection*{如果跳过本章会怎样}
\begin{itemize}
  \item \textbf{“两台机器人各自跟踪轨迹，物体被挤爆了”}：两臂抓一根钢管各自用阻抗跟踪自己末端的期望轨迹，仿真“跟得很好”，真机上钢管要么被越夹越紧（电流飙升、夹具变形）、要么两臂“打架”剧烈抖动。根因是两个独立末端阻抗\emph{没有协调内力}——期望位置差 $1$ 毫米，在刚性钢管上就放大成几百牛挤压。没有 \cref{sec:cm-internal} 的内力/外力分解，你不知道为什么、更不知道用零空间投影把内力单独拎出来控制。
  \item \textbf{“加了第三台机器人，系统突然不稳了”}：双机搬运很稳，为搬更重负载加第三台，系统开始低频振荡甚至发散，逐个调参越调越糟。根因是多机耦合系统稳定性\emph{不是}各子系统稳定性的叠加——三台通过物体形成的反馈回路可能注入净能量，破坏整体无源性。没有 \cref{sec:cm-passivity} 的无源性分析，你只能盲目试错。
  \item \textbf{“想让机器人队伍零通信协作，却不知从何下手”}：你看到“2--10 台 Go2 零通信搬运”这类工作很想复现，却完全不理解机器人之间不通信怎么知道该往哪用力。根因是没建立“被搬物体 $=$ 隐式通信信道”的认知模型。没有 \cref{sec:cm-decentral} 的去中心化力控原理，“零通信协作”对你只是一个魔法名词。
\end{itemize}

\begin{table}[H]\centering\small
\caption{本章阅读方式建议}
\label{tab:cm-reading}
\begin{tabular}{lll}
\toprule
阅读方式 & 时间 & 适合谁 \\
\midrule
精读（逐节读动机$\to$推导，亲手推抓取矩阵零空间、内力分解正交性、物体阻抗律与无源性证明） & 18--24 小时 & 第一次系统学多机协同搬运与力控 \\
速读（读结论与关键 boxed 公式、架构对比表，跳过完整证明与代码细节） & 5--7 小时 & 有单臂力控基础、想建立多机搬运全局图景 \\
速查（只看 boxed 公式 $+$ 内力分解 $+$ 架构对比 $+$ 故障排查） & 50--70 分钟 & 已学过，回查特定公式或结论 \\
\bottomrule
\end{tabular}
\end{table}

\subsection*{科研发展脉络}
多机协同搬运是\emph{四十年理论积淀（1986 起）$+$ 近五年学习革命}的领域。下表勾勒主线（论文归属、年份、卷期、数值均按源稿重述，存疑处登台账，待联网核对）。

{\small\centering
\begin{longtable}{C{0.9cm}p{4.6cm}C{2.0cm}p{4.6cm}}
\caption{协同搬运与多机力控发展脉络（按源稿重述，待核）}\label{tab:cm-history}\\
\toprule
\multicolumn{1}{c}{年份} & \multicolumn{1}{c}{论文 / 工作} & \multicolumn{1}{c}{Venue} & \multicolumn{1}{c}{核心贡献} \\
\midrule
\endfirsthead
\toprule
\multicolumn{1}{c}{年份} & \multicolumn{1}{c}{论文 / 工作} & \multicolumn{1}{c}{Venue} & \multicolumn{1}{c}{核心贡献} \\
\midrule
\endhead
\midrule
\multicolumn{4}{r}{\small\itshape （续下页）}\\
\endfoot
\bottomrule
\endlastfoot
1986 & Orin \& Oh，含闭合运动链机构的力分配控制 & ASME JDSMC & 奠基：闭链系统力分配一般公式（原文经线性规划优化能耗/负载均衡，非主打伪逆）\rebuilt{年份与卷期错：实为 1981、ASME JDSMC vol.\,103(2):134--141；"1986"系与 Oh \& Orin 1986 ICRA 闭链仿真混淆。bib 键 orin1986force 年份待改} \\
1988 & Walker, Freeman, Marcus，多机协作的内部物体加载 & ICRA & 非挤压伪逆：指出 Moore--Penrose 解会挤压物体，构造非挤压分配\pz{年份应为 1989 非 1988（同组 1988 ICRA 另有"动态任务分配"一文，易混）；"唯一"非挤压可商榷} \\
1988 & Khatib，多末端机器人系统的物体操作 & ISRR & 增广物体模型：把物体动力学纳入操作空间控制 \\
1996 & Caccavale, Chiacchio 等，双臂协调控制系列 & IEEE T-RA & 对称构型（绝对/相对运动解耦）$+$ object impedance 雏形\pz{存疑：1996 IEEE T-RA 出处不实——该组 1996 协同作实为 Chiacchio 等 ASME JDSMC，对称构型应归 Uchiyama \& Dauchez (1988/1993)；bib 实取 caccavale2008six(2008 IEEE/ASME T-Mech)} \\
2000 & Khatib 等，运动与力控统一方法 & IEEE T-RA & 操作空间内力/运动解耦，\cref{sec:cm-internal} 分解的理论母体\pz{此条即 Khatib \emph{1987} 操作空间奠基作（IEEE J. Robotics \& Automation 3(1):43--53，单作者），非"2000 IEEE T-RA"；本章正文引 khatib1987operational 键正确} \\
2008/12 & Wimb\"ock, Ott, Hirzinger，物体级抓取控制器比较 & DLR / IJRR & 物体级阻抗统一框架，多指通用，内力沿抓取线 \\
2019 & Verginis, Dimarogonas，经内力调节的协作操作：刚度理论视角 & RA-L / arXiv:1911.01297 & 抓取矩阵与刚度矩阵的 range-nullspace 关系，能量最优内力\pz{发表 venue 应为 IEEE T-CNS 10(3):1222--1233(2023)，非 RA-L；arXiv:1911.01297 正确，作者另含 Zelazo} \\
2020 & Culbertson, Schwager 等，刚体协作操作的去中心化自适应控制 & IEEE T-RO & 去中心化自适应：免负载惯性，每 agent 本地自适应（T-RO 正式年为 2021，作者另含 Slotine） \\
2018 & Wang, Schwager，经隐式通信的协作多机运输 & Frontiers Robot. AI & 隐式通信：物体运动即通信，无显式通信的力协调\rebuilt{张冠李戴：与此题/venue 对应的 Frontiers 2018 文(\emph{Collaborative Multi-Robot Transportation in Obstacle-Cluttered Environments via Implicit Communication}, frobt.2018.00090) 实为 Bechlioulis \& Kyriakopoulos；Wang \& Schwager 的对应工作是 2014 DARS 无通信操作 / 2016 IJRR Force-ANTS。键 wang2018collaborative 作者与题不自洽，待重置} \\
2021 & Carey, Werfel，经隐式通信的集体运输 & ICRA & 集体运输：leader 起头、follower 经隐式协调配合\pz{"Force-ANTS/力放大"标签张冠李戴：Force-Amplifying N-robot System 实为 Wang \& Schwager 2016 IJRR；Carey \& Werfel 2021 ICRA 真题为 \emph{Collective Transport of Unconstrained Objects via Implicit Coordination and Adaptive Compliance}，机制是自适应柔顺而非力放大} \\
2022 & Sombolestan, Nguyen，多四足对绳牵负载的协作导航与操作 & IROS & 多四足 $+$ 绳索：分布式力控搬运\rebuilt{张冠李戴：此题(\emph{...a Cable-towed Load by Multiple Quadrupedal Robots})实为 Yang 等(UC Berkeley)，发表于 RA-L 7(4):10041--10048(2022) 非 IROS，且核心是级联轨迹优化而非分布式力控；Sombolestan \& Nguyen(USC) 的相关作是 IROS 2023 刚体自适应力控（另文）。键 sombolestan2022collaborative 作者/venue/内容均待改} \\
2024 & Turrisi 等，PACC：高负载协作搬运的被动臂方法 & IROS / arXiv:2403.19862 & 被动臂 $+$ MPC：被动关节吸收耦合，简化建模\pz{被动臂为 3 转动关节(yaw-pitch-pitch)$+$弹簧阻尼、非"万向节"；各四足跑的是含臂动力学与外力估计的\emph{增广}MPC、非标准 MPC} \\
2025 & Pandit 等，decPLM：多四足协作物体运输 & arXiv & 2--10 Go2 零通信，pinch-lift-move 策略，IsaacLab 训练 \\
\end{longtable}}

\noindent\textbf{实验室脉络}（据源稿）：OSU（Orin 闭链力分配原创）$\to$ Stanford ASL（Khatib 操作空间 $+$ 增广物体）$\to$ DLR（Caccavale/Wimb\"ock/Ott 物体级阻抗）$\to$ KTH（Verginis/Dimarogonas 刚度理论 $+$ 分布式）$\to$ Stanford MSL（Schwager 去中心化自适应 $+$ 隐式通信）$\to$ Harvard（Werfel 群体搬运）$\to$ Notre Dame/USC（Sombolestan/Nguyen 多足绳索）$\to$ IIT DLS（Turrisi PACC）$\to$ 学习社区（decPLM 零通信学习）。\rebuilt{机构归属三处有误：Khatib 在 Stanford 的实验室是 Stanford Robotics Lab(SRL)、非"ASL"(ASL=Pavone 的 Autonomous Systems Lab)；Caccavale 实为意大利 Univ.\ Basilicata、非 DLR（DLR 节点应只留 Wimb\"ock/Ott）；Sombolestan/Nguyen 仅在 USC，"Notre Dame"系杜撰。其余 6 节点(OSU/KTH/Stanford MSL/Harvard/IIT DLS/学习社区)归属属实。}一条清晰主线贯穿其中：\emph{决策从“集中式精确力分配”走向“去中心化隐式协调”，建模从“精确已知负载动力学”走向“自适应/学习免模型”}——每一步都在减少对全局信息与精确模型的依赖。

% ════════════════════════════════════════════════════════════════
\section{闭链约束与抓取矩阵}
\label{sec:cm-grasp}

\paragraph{动机：两台机器人抬一根钢梁。}
设想最简单的协同搬运：两台机械臂末端各自牢牢抓住一根钢梁的两端，要把它从地上平稳抬到桌面。抓住钢梁之前，两台机械臂是\emph{独立}系统——可让左臂末端往任意方向动、右臂纹丝不动。但抓住的那一刻一切都变了：钢梁是刚体，两端被两末端“焊”住。若左臂末端上抬 $1$ 厘米、右臂保持不动，钢梁就被要求“左端上移、右端不动”，即绕右端转一个角度；若钢梁真是刚体、抓持真是刚性，右臂末端的位姿就\emph{被左臂运动和钢梁几何唯一决定}——右臂不再有“保持不动”的自由。这就是\textbf{闭合运动链（closed kinematic chain）}：左臂$\to$钢梁$\to$右臂$\to$（经共同地面）$\to$左臂构成一个环，环上任一处运动都沿环传播、约束其余各处。单臂操作是“开链”（末端自由动），多机搬运是“闭链”（所有末端被物体硬性绑定）。我们真正想控制的是\emph{物体}（钢梁该沿某轨迹运动），能直接施加的却是\emph{各末端的力}，于是两个核心问题浮现：（i）运动学——欲使物体以速度 $\mathbf{v}_o$ 运动，各末端须以什么速度运动才“兼容”物体而不撕裂闭链？（ii）静力学——各末端接触力 $\mathbf{f}_i$ 合起来对物体产生多大合力与合力矩？抓取矩阵 $\mathbf{G}$ 同时回答这两个问题，是连接“物体世界”与“机器人世界”的桥梁。

\paragraph{反面：不建立抓取矩阵会怎样。}
若固执地让每台机器人“各自为政”、独立跟踪一条预先算好的笛卡尔轨迹，会发生三件糟糕的事。\textbf{其一，过约束与内力爆炸}：一个自由刚体有 $6$ 个自由度，两个末端若都用 $6$ 维位姿独立控制，就是 $12$ 个被控量去约束 $6$ 自由度刚体——多出的 $6$ 维约束无处安放，全部变成\emph{内力}（互相挤压/拉扯/扭转钢梁的力），哪怕两条轨迹只差一个浮点误差，在刚性钢梁上也会被放大成巨大内力。\textbf{其二，无法分配负载}：“各自跟踪轨迹”里根本没有“重量分配”这个旋钮，每台只知自己的位置误差、不知物体总重与自己该出多少力。\textbf{其三，无法定义“协调”}：“协调搬运”的数学定义到底是什么？没有 $\mathbf{G}$ 你说不清，也就无法检验两台是否协调。这三件事都指向同一个缺失的工具——抓取矩阵。

\begin{insight}[抓取矩阵是闭链约束的数学化身]
抓取矩阵不是“为方便计算而引入的工具矩阵”，而是闭链约束的\emph{数学化身}。$\dim\ker\mathbf{G}=3N-6$ 这个数字——接触力空间维度 $3N$ 减物体自由度 $6$——精确量化了“过约束的程度”：有多少维的力可以施加却不影响物体运动。这 $3N-6$ 维就是内力的容身之所，也是搬运一切麻烦的根源与一切精巧控制的着力点。
\end{insight}

\paragraph{历史：两条源流汇成统一框架。}
抓取矩阵有两条历史源流。一条来自\emph{腿足机器人与闭链机构}：Orin 与 Oh\cite{orin1986force} 研究含闭合运动链的机器人（如四足站立时四腿 $+$ 地面 $+$ 躯干构成的闭链）时，提出闭链系统力分配的一般公式，最早系统处理“接触力多于任务自由度”导致的欠定问题。\rebuilt{该文年份应为 1981（非源稿"1986"），ASME JDSMC vol.\,103(2):134--141；且其求解用\emph{线性规划}优化能耗/负载均衡，"用伪逆求解"的描述不准——伪逆/一般解处理更属后续 Orin 学派工作（如 Cheng \& Orin 1991, IEEE T-SMC 21(1):25--32）。}另一条来自\emph{多指抓取与多臂操作}：Walker、Freeman、Marcus\cite{walker1988internal} 在 ICRA 指出用 Moore--Penrose 伪逆分配会在物体内部产生不必要挤压（squeeze），并构造“非挤压”力分配\pz{年份应为 1989 非 1988（故"1988 是关键一年"措辞欠准；同组 1988 ICRA 另有"动态任务分配"一文易混）；"唯一"非挤压属原文主张、后续工作(Chung 等)有异议，宜软化}；Khatib\cite{khatib1988object} 在 ISRR(1988) 提出“增广物体（augmented object）”模型，把被操作物体动力学纳入操作空间控制\cite{khatib1987operational}。两条线在 1990 年代被 Caccavale、Chiacchio 等\cite{caccavale2008six} 统一到双臂协调控制，“抓取矩阵”作为标准术语固定下来。

\begin{note}[类比的边界：多指抓取 vs 多臂搬运]
多指手抓取（grasping）与多臂搬运（cooperative manipulation）共享同一套抓取矩阵数学，这正是本章与双臂力控用同一个 $\mathbf{G}$ 的原因。\emph{像的部分}：都是多接触点对同一刚体施力，都用 $\mathbf{w}=\mathbf{G}\mathbf{f}$ 描述力映射，内力都在 $\ker\mathbf{G}$。\emph{不像的部分}：多指抓取的接触通常是\textbf{单边的}（手指只能推不能拉，受摩擦锥约束，\cref{sec:cm-qp}），而多臂刚性抓持是\textbf{双边的}（夹爪焊死物体，可推可拉可扭）。不要把“刚性抓持下内力可任意”延伸到“指尖接触”——指尖接触的内力必须保持在摩擦锥内，否则物体滑落。
\end{note}

\subsection{静力学：接触力合成物体旋量}
\label{sec:cm-grasp-statics}
考虑被搬物体 $\mathcal{O}$ 为刚体，质心在 $\mathbf{p}_o$。第 $i$ 个机器人在接触点 $\mathbf{x}_i$ 处对物体施加接触力 $\mathbf{f}_i\in\mathbb{R}^3$（先只考虑力、不考虑接触力矩，对应“点接触/球铰抓持”模型；能传力矩的“刚性抓持”留作 \cref{sec:cm-grasp-rigid} 的推广）。该接触力对物体质心产生两个效应：\emph{平移效应}——力 $\mathbf{f}_i$ 直接贡献合力；\emph{转动效应}——力作用在偏离质心的 $\mathbf{x}_i$ 上，产生力矩 $\boldsymbol{\tau}_i=\mathbf{r}_i\times\mathbf{f}_i$，其中 $\mathbf{r}_i=\mathbf{x}_i-\mathbf{p}_o$。把 $N$ 个接触点贡献叠加，作用在质心的合力旋量 $\mathbf{w}_o=[\mathbf{F}_o^\top,\boldsymbol{\tau}_o^\top]^\top$ 为
\begin{equation}
  \mathbf{F}_o=\sum_{i=1}^N\mathbf{f}_i,\qquad \boldsymbol{\tau}_o=\sum_{i=1}^N\mathbf{r}_i\times\mathbf{f}_i.
  \label{eq:cm-wrench-sum}
\end{equation}
把叉积写成矩阵乘法。叉积 $\mathbf{r}_i\times(\cdot)$ 对第二参数是线性算子，可写成反对称矩阵 $\skewm{\mathbf{r}}_i$ 左乘（与 \cref{ch:lie} 的 $(\cdot)^\wedge$ 同义）：
\begin{equation}
  \skewm{\mathbf{r}}_i=
  \begin{bmatrix}0&-r_{i,z}&r_{i,y}\\ r_{i,z}&0&-r_{i,x}\\ -r_{i,y}&r_{i,x}&0\end{bmatrix},
  \qquad \skewm{\mathbf{r}}_i\mathbf{f}_i=\mathbf{r}_i\times\mathbf{f}_i.
  \label{eq:cm-skew}
\end{equation}
反对称（$\skewm{\mathbf{r}}_i^\top=-\skewm{\mathbf{r}}_i$）来自叉积的反交换律 $\mathbf{r}\times\mathbf{f}=-\mathbf{f}\times\mathbf{r}$；三列由 $\mathbf{r}\times\mathbf{e}_j$ 求得。于是合力旋量写成堆叠形式：
\begin{equation}
\boxed{\;
  \mathbf{w}_o=\begin{bmatrix}\mathbf{F}_o\\ \boldsymbol{\tau}_o\end{bmatrix}
  =\underbrace{\begin{bmatrix}\mathbf{I}_3&\mathbf{I}_3&\cdots&\mathbf{I}_3\\ \skewm{\mathbf{r}}_1&\skewm{\mathbf{r}}_2&\cdots&\skewm{\mathbf{r}}_N\end{bmatrix}}_{\mathbf{G}\,\in\,\mathbb{R}^{6\times 3N}}
  \underbrace{\begin{bmatrix}\mathbf{f}_1\\ \mathbf{f}_2\\ \vdots\\ \mathbf{f}_N\end{bmatrix}}_{\mathbf{f}\,\in\,\mathbb{R}^{3N}}=\mathbf{G}\mathbf{f}.
\;}
  \label{eq:cm-grasp}
\end{equation}
这就是\textbf{抓取矩阵方程}。$\mathbf{G}$ 的第 $i$ 个 $6\times 3$ 分块是
\begin{equation}
  \mathbf{G}_i=\begin{bmatrix}\mathbf{I}_3\\ \skewm{\mathbf{r}}_i\end{bmatrix}\in\mathbb{R}^{6\times 3},
  \label{eq:cm-Gi}
\end{equation}
上半块 $\mathbf{I}_3$ 把接触力原样搬到合力（“力就是力”），下半块 $\skewm{\mathbf{r}}_i$ 把接触力经力臂转成力矩（“杠杆效应”）。务必看清：$\mathbf{G}$ \emph{只依赖几何}（各接触点相对质心的 $\mathbf{r}_i$），不含任何动力学参数（物体质量、惯量都没出现）——这是 $\mathbf{G}$ 作为“纯运动学/几何”对象的本质。

\subsection{运动学对偶：物体速度分发到接触点速度}
\label{sec:cm-grasp-kin}
反过来看运动。物体作为刚体，质心以广义速度 $\mathbf{v}_o=[\dot{\mathbf{p}}_o^\top,\boldsymbol{\omega}_o^\top]^\top$ 运动。刚体上任一点 $\mathbf{x}_i$ 的速度由刚体运动学给出：
\begin{equation}
  \dot{\mathbf{x}}_i=\dot{\mathbf{p}}_o+\boldsymbol{\omega}_o\times\mathbf{r}_i=\dot{\mathbf{p}}_o-\mathbf{r}_i\times\boldsymbol{\omega}_o=\dot{\mathbf{p}}_o-\skewm{\mathbf{r}}_i\boldsymbol{\omega}_o.
  \label{eq:cm-pt-vel}
\end{equation}
写成矩阵形式（注意 $\boldsymbol{\omega}_o\times\mathbf{r}_i=-\mathbf{r}_i\times\boldsymbol{\omega}_o=-\skewm{\mathbf{r}}_i\boldsymbol{\omega}_o$，符号要小心）：
\begin{equation}
  \dot{\mathbf{x}}_i=\begin{bmatrix}\mathbf{I}_3&-\skewm{\mathbf{r}}_i\end{bmatrix}\begin{bmatrix}\dot{\mathbf{p}}_o\\ \boldsymbol{\omega}_o\end{bmatrix}
  =\begin{bmatrix}\mathbf{I}_3&\skewm{\mathbf{r}}_i^\top\end{bmatrix}\mathbf{v}_o=\mathbf{G}_i^\top\mathbf{v}_o,
  \label{eq:cm-vel-dual-step}
\end{equation}
末步用 $\skewm{\mathbf{r}}_i^\top=-\skewm{\mathbf{r}}_i$。所以接触点速度与物体速度的关系是
\begin{equation}
\boxed{\;\dot{\mathbf{x}}_i=\mathbf{G}_i^\top\mathbf{v}_o,\qquad \dot{\mathbf{X}}=\mathbf{G}^\top\mathbf{v}_o,\;}
  \label{eq:cm-vel-dual}
\end{equation}
其中 $\dot{\mathbf{X}}=[\dot{\mathbf{x}}_1^\top,\dots,\dot{\mathbf{x}}_N^\top]^\top$。\textbf{这就是闭链约束的运动学形式}：所有接触点速度都必须是同一个物体速度 $\mathbf{v}_o$ 经 $\mathbf{G}^\top$ 的像。若各末端速度 $\dot{\mathbf{x}}_i$ 不满足“共同来自某个 $\mathbf{v}_o$”，它们就在试图撕裂刚体——物理上不可能，结果就是内力。\cref{eq:cm-vel-dual} 是“协调”的精确数学定义。

\subsection{对偶性：为何力映射与速度映射由同一个 \texorpdfstring{$\mathbf{G}$}{G} 联系}
\label{sec:cm-grasp-duality}
\cref{eq:cm-grasp} 用 $\mathbf{G}$、\cref{eq:cm-vel-dual} 用 $\mathbf{G}^\top$，这不是巧合，而是\textbf{功的守恒（运动静力学对偶，kineto-statics duality）}。考虑接触界面传递的瞬时功率：从接触力一侧看是 $\sum_i\mathbf{f}_i^\top\dot{\mathbf{x}}_i=\mathbf{f}^\top\dot{\mathbf{X}}$；从物体一侧看是 $\mathbf{w}_o^\top\mathbf{v}_o$。刚性无滑动接触不储存也不耗散能量，两者必须相等：
\begin{equation}
  \mathbf{f}^\top\dot{\mathbf{X}}=\mathbf{w}_o^\top\mathbf{v}_o
  \;\Longrightarrow\;
  \mathbf{f}^\top(\mathbf{G}^\top\mathbf{v}_o)=(\mathbf{G}\mathbf{f})^\top\mathbf{v}_o
  \;\Longrightarrow\;
  \mathbf{f}^\top\mathbf{G}^\top\mathbf{v}_o=\mathbf{f}^\top\mathbf{G}^\top\mathbf{v}_o.\ \checkmark
  \label{eq:cm-duality}
\end{equation}
恒等式成立——这证明：\emph{只要力映射是 $\mathbf{G}$，速度映射就必然是 $\mathbf{G}^\top$}。这与操作空间动力学里 $\boldsymbol{\tau}=\mathbf{J}^\top\mathbf{F}$（力）和 $\dot{\mathbf{x}}=\mathbf{J}\dot{\mathbf{q}}$（速度）由同一个 $\mathbf{J}$ 联系是完全相同的对偶结构。

\begin{insight}[对偶性的统一视角]
单臂的 $\mathbf{J}$、多机的 $\mathbf{G}$、多指手的抓取映射，本质都是\emph{一个把“低维任务空间”和“高维执行器空间”联系起来的线性映射，其转置自动给出对偶方向的映射}。记住这个统一视角，你就不必为每个新场景重新推导：看到力映射 $\mathbf{A}$，立刻知道速度映射是 $\mathbf{A}^\top$；看到执行器空间维度 $n$ 大于任务空间维度 $m$，立刻知道零空间维度是 $n-m$，那里住着“不影响任务的执行器自由度”——单臂叫冗余自运动，多机叫内力。
\end{insight}

\subsection{零空间维度：为何恒等于 \texorpdfstring{$3N-6$}{3N-6}}
\label{sec:cm-grasp-null}
抓取矩阵 $\mathbf{G}\in\mathbb{R}^{6\times 3N}$。当 $N\ge 2$ 时 $3N\ge 6$，$\mathbf{G}$ 是“宽矩阵”（列多于行）。只要抓取点不退化（不全部共线、不全部重合），$\mathbf{G}$ 行满秩，$\rank(\mathbf{G})=6$。由秩-零化度定理（rank--nullity theorem）：
\begin{equation}
  \dim\ker\mathbf{G}=3N-\rank(\mathbf{G})=3N-6.
  \label{eq:cm-null-dim}
\end{equation}
这个数字是本章的灵魂：在 $3N$ 维接触力空间里，有 $6$ 维力“穿透” $\mathbf{G}$ 变成物体运动（外力/运动力），剩下 $3N-6$ 维力被 $\mathbf{G}$ 映射为零——施加它们物体一动不动，它们全部是\textbf{内力}。下表举具体数字建立直觉。

\begin{table}[H]\centering\small
\caption{点接触模型下内力维度随机器人数的变化}
\label{tab:cm-null-examples}
\begin{tabular}{lcccc}
\toprule
场景 & $N$ & 接触力空间 $3N$ & 物体自由度 & 内力维度 $3N-6$ \\
\midrule
双臂搬运（点接触） & 2 & 6 & 6 & \textbf{0} \\
三机器人搬运 & 3 & 9 & 6 & \textbf{3} \\
四足搬运 / 四指抓取 & 4 & 12 & 6 & \textbf{6} \\
五机器人搬运 & 5 & 15 & 6 & \textbf{9} \\
\bottomrule
\end{tabular}
\end{table}

\noindent 注意一个反直觉结论：\emph{两台机器人用点接触（每点 $3$ 维力）抓物体，内力维度是 $0$}——没有内力！这看似与“双臂搬运会挤爆物体”矛盾。关键在“点接触”假设：若每个接触只能传 $3$ 维力（不能传力矩），$N=2$ 时 $3N-6=0$ 确实无内力可调。但真实双臂用夹爪是\textbf{刚性抓持}，每个接触能传 $6$ 维力旋量（力 $+$ 力矩），此时 $\mathbf{G}\in\mathbb{R}^{6\times 6N}$，内力维度变成 $6N-6$，$N=2$ 时是 $6$ 维——这才是双臂搬运产生内力的来源（见 \cref{sec:cm-grasp-rigid}）。

\subsection{推广：从点接触到刚性抓持}
\label{sec:cm-grasp-rigid}
真实搬运中机械臂夹爪牢牢夹住物体属于\textbf{刚性抓持}——接触界面不仅传力，还传\emph{力矩}。此时第 $i$ 个接触传递 $6$ 维力旋量 $\mathbf{w}_i=[\mathbf{f}_i^\top,\mathbf{m}_i^\top]^\top\in\mathbb{R}^6$，抓取矩阵变为 $\mathbf{G}\in\mathbb{R}^{6\times 6N}$，第 $i$ 块是
\begin{equation}
  \mathbf{G}_i^{\text{rigid}}=\begin{bmatrix}\mathbf{I}_3&\mathbf{0}\\ \skewm{\mathbf{r}}_i&\mathbf{I}_3\end{bmatrix}\in\mathbb{R}^{6\times 6},
  \label{eq:cm-Gi-rigid}
\end{equation}
下半块多了一个 $\mathbf{I}_3$：接触力矩 $\mathbf{m}_i$ 直接贡献到物体力矩（不经力臂）。此时内力维度变为 $6N-6$。下表把“接触模型 $\to$ 单点力维度 $\to$ 抓取矩阵尺寸 $\to$ 内力维度”系统化（穷举式分类，而非随意举例）。

\begin{table}[H]\centering\small
\caption{接触模型与内力维度的系统分类}
\label{tab:cm-contact-models}
\begin{tabular}{p{3.0cm}p{3.2cm}cccp{2.6cm}}
\toprule
接触模型 & 物理 & 单点力维 $c$ & $\mathbf{G}$ 尺寸 & 内力维度（$N{=}2$） \\
\midrule
点接触无摩擦 & 只传法向力 & 1 & $6\times N$ & $N-6$（$N{<}6$ 不能完全约束） \\
点接触带摩擦（PCWF） & 传 $3$ 维力，受摩擦锥约束 & 3 & $6\times 3N$ & $3N-6$（$N{=}2$：$0$） \\
软指接触（SF） & 传 $3$ 维力 $+$ 法向扭矩 & 4 & $6\times 4N$ & $4N-6$（$N{=}2$：$2$） \\
刚性抓持 & 传 $6$ 维力旋量 & 6 & $6\times 6N$ & $6N-6$（$N{=}2$：$6$） \\
\bottomrule
\end{tabular}
\end{table}

\noindent 本章后续除非特别说明，以\textbf{点接触带摩擦（每点 $3$ 维力，$\mathbf{G}\in\mathbb{R}^{6\times 3N}$）}为默认模型（多足搬运、绳索搬运、末端球铰抓持都属此类）；涉及双臂刚性抓持时切换到 $6$ 维模型。双臂协调力控用的就是刚性抓持版本，与此处推广一致。统一记内力维度为 $cN-6$。

\paragraph{代码实现：抓取矩阵的组装。}
理论讲完，用代码确认。维度运行期才知道（$N$ 可变）用动态矩阵，反对称块用固定大小（栈分配）。核心是组装 \cref{eq:cm-grasp} 的 $\mathbf{G}$，并写一个反对称矩阵函数 \verb@skew(r)@ 满足 \verb@skew(r)*v == r.cross(v)@：每帧用当前接触点与质心重算 $\mathbf{r}_i=$ 接触点 $-$ 质心（\textbf{方向勿写反}），上块填 $\mathbf{I}_3$、下块填 \verb@skew(r_i)@。三处常见错误：（1）把 $\mathbf{r}_i$ 写成 质心 $-$ 接触点（符号反 $\Rightarrow$ 力矩贡献整体反号，平移搬运正常、一旋转就发散，因为力对、只有力矩错）；（2）物体转动后接触点相对质心几何在世界系会变，若缓存初始 $\mathbf{r}_i$ 不更新则 $\mathbf{G}$ 过时——应每帧重算，或在体坐标系固定 $\mathbf{r}_i^{\text{body}}$、用当前姿态 $\mathbf{R}_o$ 旋转得 $\mathbf{r}_i=\mathbf{R}_o\mathbf{r}_i^{\text{body}}$；（3）反对称矩阵未一次性初始化全 $9$ 元素（含 $3$ 个 $0$），漏项留栈垃圾值。代码占比按规范从略，要点已述。

\begin{insight}[抓取矩阵的几何决定工程性能]
正因 $\mathbf{G}$ 只依赖接触点几何 $\mathbf{r}_i$，工程上\emph{接触点的几何布局直接决定 $\mathbf{G}$ 的条件数与内力空间的“形状”}。若四台四足把附着点布得太靠近（$\mathbf{r}_i$ 都很小且方向接近），$\mathbf{G}$ 的下半块 $\skewm{\mathbf{r}}_i$ 都很“弱”，物体转动控制权威性低、$\mathbf{G}\mathbf{G}^\top$ 病态，力分配（\cref{sec:cm-qp}）会放大噪声。所以多四足绳索搬运、decPLM 的 pinch-lift-move 都隐含要求附着点“撑开”以获得良好力矩臂——这正是 $\mathbf{G}$ 的几何在工程中的回响。
\end{insight}

\begin{pitfall}[编程陷阱：反对称矩阵符号反号导致“一转就翻”]
把 $\mathbf{r}_i$ 写成 质心 $-$ 接触点，或 $\skewm{\mathbf{r}}$ 的上/下三角符号写反（常见：把 $-r_z$ 写成 $r_z$），会导致：纯平移搬运（只用力、不用力矩）一切正常，但只要物体需要旋转，力矩方向整体反号，控制器把物体往反方向越推越远，表现为“一旋转就发散/翻转”。因为力（上块 $\mathbf{I}_3$）是对的、bug 藏在力矩（下块）里，平移测试根本暴露不出来。根因：下半块 $\skewm{\mathbf{r}}_i$ 编码“力臂叉乘”，叉积反交换，$\mathbf{r}_i$ 反号或 $\skewm{\mathbf{r}}$ 符号反号都让 $\boldsymbol{\tau}=\mathbf{r}\times\mathbf{f}$ 整体变 $-\boldsymbol{\tau}$。对策：严格用 $\mathbf{r}_i=$ 接触点 $-$ 质心；单元测试验证 $\skewm{\mathbf{r}}\mathbf{v}=\mathbf{r}\times\mathbf{v}$ 对随机 $\mathbf{r},\mathbf{v}$ 成立。
\end{pitfall}

\begin{pitfall}[概念误区：以为“两台机器人点接触搬运一定有内力”]
新手常想“双臂搬钢管会挤爆物体 $\Rightarrow$ 双臂搬运总有内力要控制”。实际上内力维度 $=cN-6$：两台机器人若是点接触（$c{=}3$），内力维度 $=3\cdot2-6=0$，数学上根本没有内力自由度！“挤爆物体”发生在刚性抓持（$c{=}6$，内力维度 $6$）或两端用夹爪传力矩的情形。搞错接触模型会让你为一个不存在的内力空间设计控制器（点接触双机），或漏掉真实存在的内力（刚性抓持双机）。延伸：三台机器人点接触（$c{=}3$）内力维度 $=3$，此时即便点接触也有内力（三台可互相“对拉”而不动物体）——“几台开始有内力”取决于 $cN$ 何时超过 $6$。
\end{pitfall}

\begin{pitfall}[思维陷阱：把抓取矩阵当成“又一个要背的矩阵”]
孤立记忆 $\mathbf{G}=[\mathbf{I};\skewm{\mathbf{r}}]$ 而不和雅可比建立联系，是常见误区。$\mathbf{G}$ 和单臂雅可比 $\mathbf{J}$ 是同一类对象——任务空间$\leftrightarrow$执行器空间的线性映射，转置给对偶方向：$\mathbf{J}:\dot{\mathbf{q}}\to\dot{\mathbf{x}}$（速度）/ $\boldsymbol{\tau}=\mathbf{J}^\top\mathbf{F}$（力）；$\mathbf{G}:\mathbf{f}\to\mathbf{w}$（力）/ $\dot{\mathbf{X}}=\mathbf{G}^\top\mathbf{v}_o$（速度）。学新映射时问三件事：（1）力往哪个方向映射（确定 $\mathbf{G}$ 还是 $\mathbf{G}^\top$）？（2）哪边维度高（确定零空间在哪边）？（3）零空间的物理意义是什么？而非死记公式。自检：不看书，仅凭“功守恒”从 $\mathbf{w}=\mathbf{G}\mathbf{f}$ 推出 $\dot{\mathbf{X}}=\mathbf{G}^\top\mathbf{v}_o$；仅凭秩-零化度说出任意接触模型的内力维度。
\end{pitfall}

\begin{practice}
\begin{enumerate}
  \item （$\star\star$）三台机器人点接触抓持刚体，接触点相对质心分别为 $\mathbf{r}_1=(1,0,0)$、$\mathbf{r}_2=(-0.5,0.866,0)$、$\mathbf{r}_3=(-0.5,-0.866,0)$（正三角形，单位米）。手写 $9\times 3$ 的 $\mathbf{G}^\top$（即 $\dot{\mathbf{X}}=\mathbf{G}^\top\mathbf{v}_o$），验证 $\rank(\mathbf{G})=6$、$\dim\ker\mathbf{G}=3$，并给出 $\ker\mathbf{G}$ 一组基向量的物理解释（提示：三台“沿三条边互相对拉”且合力合力矩为零的力模式）。
  \item （$\star\star$）把抓取矩阵组装扩展为刚性抓持版本（每点 $6$ 维力旋量，$\mathbf{G}\in\mathbb{R}^{6\times 6N}$）。随机生成 $N=2$ 的接触点与一组接触力旋量，验证（a）$\mathbf{G}\mathbf{f}$ 的力部分等于各 $\mathbf{f}_i$ 之和；（b）施加一对“等大反向、沿连线”的力（纯内力）时 $\mathbf{G}\mathbf{f}=\mathbf{0}$。报告你构造的纯内力向量。
  \item （$\star\star\star$）物体运动中姿态不断变化。设计在物体体坐标系固定存储 $\mathbf{r}_i^{\text{body}}$、每帧输入当前姿态 $\mathbf{R}_o$ 输出世界系 $\mathbf{G}$ 的方案。讨论：为何体坐标系存储比每帧重测世界系接触点更鲁棒？若某台机器人搬运中打滑、接触点偏移，模型会出什么问题、如何检测？
\end{enumerate}
\end{practice}

% ════════════════════════════════════════════════════════════════
\section{内力与外力分解}
\label{sec:cm-internal}

\paragraph{动机：同一个搬运，无穷多种用力方式。}
回到两台机器人抬钢梁。设钢梁重 $100$ 牛，要匀速向上抬（需向上 $100$ 牛合力抵消重力）。下面三种用力方式都能完成任务：\emph{方式甲}两台各向上出 $50$ 牛（合力 $100$ 牛向上）；\emph{方式乙}左出 $80$ 牛、右出 $20$ 牛（合力仍 $100$ 牛）；\emph{方式丙}两台各向上 $50$ 牛，\emph{同时}各自向内（朝对方）水平推 $200$ 牛（竖直合力 $100$ 牛、水平合力 $0$）。三种方式对物体的\emph{运动效果完全相同}，但对物体的\emph{内部受力天差地别}：方式丙里钢梁被两侧各 $200$ 牛水平挤压——若钢梁是易碎玻璃管，方式丙会把它捏碎，方式甲不会。这就引出核心问题：接触力 $\mathbf{f}$ 里，哪部分决定运动、哪部分决定挤压？我们需要一个数学手术，把 $\mathbf{f}$ 精确切成“运动力”和“内力”两块分别控制——运动力由任务（物体该怎么动）决定，内力由抓持需求（夹多紧才不滑落、又不夹碎）决定，这两个目标正交，应独立设定。

\paragraph{反面：不区分内外力会怎样。}
若拒绝分解、直接让控制器“算出一组能实现物体运动的接触力”就用，会遇三个无法回避的问题。\textbf{其一，内力失控、物体被挤爆或脱手}：不分解就没有“内力”这个旋钮，控制器随便挑一组满足 $\mathbf{G}\mathbf{f}=\mathbf{w}$ 的解——可能恰好带巨大挤压（玻璃管碎），也可能内力为负（机器人互相“对拉”，物体被扯松甚至脱手）。\textbf{其二，解的不唯一性无法消解}：$\mathbf{G}\mathbf{f}=\mathbf{w}$ 有 $3N-6$ 维解空间，不分解就没有原则挑哪一个，“随便挑”等于“行为不可预测”。\textbf{其三，无法把“运动控制”和“抓持控制”解耦设计}：运动控制器关心“物体跟没跟上轨迹”、抓持控制器关心“夹得紧不紧”，不分解二者揉在一组力里互相干扰；分解后落在正交子空间，各调各的、互不影响。

\begin{insight}[内力/外力分解 $=$ 力空间的正交分解]
内力/外力分解的深层意义，是把一个 $3N$ 维“混沌力空间”\emph{正交分解}成两个物理意义清晰、可独立指定的子空间：
\[
  \mathbb{R}^{3N}=\underbrace{\range(\mathbf{G}^{+})}_{\text{运动力，}6\text{ 维}}\ \oplus\ \underbrace{\ker(\mathbf{G})}_{\text{内力，}(3N-6)\text{ 维}}.
\]
运动力子空间由“物体该怎么动”驱动，内力子空间由“该怎么夹”驱动。这个正交分解不是数学游戏——它是“运动”与“抓持”两个物理任务在力空间里的天然边界。
\end{insight}

\paragraph{历史：从伪逆解到非挤压分配。}
最朴素的分配是 Moore--Penrose 伪逆 $\mathbf{f}=\mathbf{G}^{+}\mathbf{w}$，给出满足 $\mathbf{G}\mathbf{f}=\mathbf{w}$ 的\emph{最小范数}解。Orin--Oh\cite{orin1986force} 的闭链力分配从此起步\rebuilt{Orin--Oh 年份应为 1981 非 1986（详见前节脉络表注）}。Walker、Freeman、Marcus\cite{walker1988internal} 在 ICRA 指出一个微妙而重要的问题：\textbf{Moore--Penrose 伪逆给出的解，一般会在物体内部产生不必要的挤压}——它“最小范数”是数学意义上的最小，并不“非挤压”；他们构造了考虑抓取几何的特殊伪逆（后称 Walker 非挤压伪逆）使内力为零\pz{该文年份应为 1989；"存在唯一的力分配使内力为零"为原文主张，后续工作对"唯一/确无内力"有异议，宜读作其原始论断}。这第一次把“内力”从“伪逆解的副产品”提升为“需要显式设计的独立量”。现代多机搬运（Caccavale 的 object impedance\cite{caccavale2008six}、Wimb\"ock--Ott 的物体级抓取控制\cite{wimbock2012comparison}、Verginis 的内力调节\cite{verginis2019cooperative}）全部建立在“内力可显式指定、且应独立于运动力设计”这一认识上。Verginis \& Dimarogonas（及 Zelazo）\cite{verginis2019cooperative} 更从\emph{刚度理论（rigidity theory）}视角揭示抓取矩阵零空间与图刚度矩阵的 range-nullspace 关系，给出能量最优内力分配\pz{arXiv:1911.01297 正确，但正式发表于 IEEE T-CNS 10(3):1222--1233(2023)，非源稿所记 RA-L}。

\subsection{分解公式}
\label{sec:cm-internal-formula}
接触力 $\mathbf{f}\in\mathbb{R}^{3N}$ 要满足 $\mathbf{G}\mathbf{f}=\mathbf{w}_{\text{ext}}$，其中 $\mathbf{w}_{\text{ext}}$ 是期望施加给物体的外部力旋量（由运动控制器算出，见 \cref{sec:cm-impedance}）。这是欠定线性方程，通解 $=$ 特解 $+$ 齐次解：
\begin{equation}
\boxed{\;
  \mathbf{f}=\underbrace{\mathbf{G}^{+}\mathbf{w}_{\text{ext}}}_{\mathbf{f}_{\text{ext}}\,(\text{运动力})}
  +\underbrace{(\mathbf{I}-\mathbf{G}^{+}\mathbf{G})\,\boldsymbol{\eta}}_{\mathbf{f}_{\text{int}}\,(\text{内力})},
  \qquad \boldsymbol{\eta}\in\mathbb{R}^{3N}\ \text{任取}.
\;}
  \label{eq:cm-decomp}
\end{equation}
\begin{itemize}
  \item \textbf{运动力（外力）} $\mathbf{f}_{\text{ext}}=\mathbf{G}^{+}\mathbf{w}_{\text{ext}}$：满足 $\mathbf{G}\mathbf{f}_{\text{ext}}=\mathbf{w}_{\text{ext}}$（确实产生期望旋量），是最小范数特解，位于 $\range(\mathbf{G}^{+})=\range(\mathbf{G}^\top)=\rowsp(\mathbf{G})$。
  \item \textbf{内力} $\mathbf{f}_{\text{int}}=(\mathbf{I}-\mathbf{G}^{+}\mathbf{G})\boldsymbol{\eta}=\mathbf{P}_N\boldsymbol{\eta}$：被 $\mathbf{G}$ 映为零（$\mathbf{G}\mathbf{f}_{\text{int}}=\mathbf{0}$，因 $\mathbf{G}\mathbf{P}_N=\mathbf{G}-\mathbf{G}\mathbf{G}^{+}\mathbf{G}=\mathbf{G}-\mathbf{G}=\mathbf{0}$），不影响物体运动，位于 $\ker\mathbf{G}$。$\boldsymbol{\eta}$ 是设计者指定的内力意图（“想夹多紧”）。
\end{itemize}
记零空间投影矩阵 $\mathbf{P}_N:=\mathbf{I}-\mathbf{G}^{+}\mathbf{G}$。当 $\mathbf{G}^{+}$ 取 Moore--Penrose 伪逆时它是到 $\ker\mathbf{G}$ 的\textbf{正交投影}，满足两个关键性质：
\begin{equation}
  \mathbf{P}_N^2=\mathbf{P}_N\ (\text{幂等}),\qquad \mathbf{P}_N^\top=\mathbf{P}_N\ (\text{对称/正交投影}).
  \label{eq:cm-PN-props}
\end{equation}

\begin{insight}[为什么 $\mathbf{f}_{\text{ext}}$ 与 $\mathbf{f}_{\text{int}}$ 正交]
$\mathbf{f}_{\text{ext}}\in\rowsp(\mathbf{G})$、$\mathbf{f}_{\text{int}}\in\ker(\mathbf{G})$，而线性代数基本定理告诉我们 $\rowsp(\mathbf{G})\perp\ker(\mathbf{G})$（行空间与零空间是正交补），所以 $\mathbf{f}_{\text{ext}}^\top\mathbf{f}_{\text{int}}=0$ 恒成立。这个正交性是“运动力与内力互不干扰”的数学根基——你改变内力意图 $\boldsymbol{\eta}$，$\mathbf{f}_{\text{int}}$ 在零空间里变，但 $\mathbf{f}_{\text{ext}}$ 纹丝不动，物体运动完全不受影响。这与单臂零空间投影“次要任务不破坏主任务”是同一个数学。
\end{insight}

\subsection{内力为何“不影响运动”——再确认}
\label{sec:cm-internal-verify}
验证 $\mathbf{f}_{\text{int}}$ 真的不产生物体旋量：
\begin{equation}
  \mathbf{G}\mathbf{f}_{\text{int}}=\mathbf{G}(\mathbf{I}-\mathbf{G}^{+}\mathbf{G})\boldsymbol{\eta}=(\mathbf{G}-\mathbf{G}\mathbf{G}^{+}\mathbf{G})\boldsymbol{\eta}.
  \label{eq:cm-int-zero}
\end{equation}
对 Moore--Penrose 伪逆，$\mathbf{G}\mathbf{G}^{+}\mathbf{G}=\mathbf{G}$（伪逆四条件之一），所以 $\mathbf{G}\mathbf{f}_{\text{int}}=(\mathbf{G}-\mathbf{G})\boldsymbol{\eta}=\mathbf{0}$。$\checkmark$ 物理上：内力是各接触点之间“自相平衡”的力——合力为零、合力矩也为零，对物体质心不产生任何净效应，只在物体内部造成应力（挤压或张紧）。

\subsection{内力子空间的基：参数化 \texorpdfstring{$(3N-6)$}{(3N-6)} 维内力}
\label{sec:cm-internal-basis}
$\mathbf{f}_{\text{int}}=\mathbf{P}_N\boldsymbol{\eta}$ 用 $3N$ 维的 $\boldsymbol{\eta}$ 参数化一个只有 $3N-6$ 维的子空间，是冗余的（$\boldsymbol{\eta}$ 里有 $6$ 维被 $\mathbf{P}_N$ 投掉）。更经济的做法是找到 $\ker\mathbf{G}$ 的一组基 $\mathbf{E}\in\mathbb{R}^{3N\times(3N-6)}$（列向量张成 $\ker\mathbf{G}$，即 $\mathbf{G}\mathbf{E}=\mathbf{0}$ 且 $\mathbf{E}$ 列满秩），用 $(3N-6)$ 维参数 $\boldsymbol{\lambda}$ 表示内力：
\begin{equation}
  \mathbf{f}_{\text{int}}=\mathbf{E}\,\boldsymbol{\lambda},\qquad \boldsymbol{\lambda}\in\mathbb{R}^{3N-6}.
  \label{eq:cm-int-basis}
\end{equation}
$\boldsymbol{\lambda}$ 的每个分量直接对应一个独立“内力模式”的强度，物理意义清晰。数值上对 $\mathbf{G}$ 做 SVD 或 QR，零奇异值对应的右奇异向量就是 $\ker\mathbf{G}$ 的标准正交基。

\noindent\textbf{三机器人正三角形的内力基（具体例子）}：三台点接触（$N=3$，内力维度 $3N-6=3$），$\ker\mathbf{G}$ 的三个基模式有清晰物理图像（见下表），每个模式合力合力矩都为零（否则就不在 $\ker\mathbf{G}$ 里），但都在物体内部产生真实应力。控制器通过设定 $\boldsymbol{\lambda}$ 选择“夹多紧”——通常取正值的径向夹紧模式，确保接触不滑脱（满足摩擦锥，\cref{sec:cm-qp}）。

\begin{table}[H]\centering\small
\caption{三机器人正三角形的三个内力模式}
\label{tab:cm-int-modes}
\begin{tabular}{llcc}
\toprule
内力模式 & 物理图像 & 合力 & 合力矩 \\
\midrule
模式 1 & 三台沿“指向质心”方向同时内推（径向夹紧） & $\mathbf{0}$ & $\mathbf{0}$ \\
模式 2 & 三台沿切向同时施力（绕质心扭转的对抗） & $\mathbf{0}$ & $\mathbf{0}$ \\
模式 3 & 一对机器人对拉、第三台平衡（边内力） & $\mathbf{0}$ & $\mathbf{0}$ \\
\bottomrule
\end{tabular}
\end{table}

\begin{note}[类比的边界：内力 vs 螺栓预紧力]
内力之于多机搬运，恰如“预紧力（preload）”之于螺栓连接。\emph{像的部分}：都是不产生宏观运动、只在结构内部维持的力，都用来保证连接/接触不松脱。\emph{不像的部分}：螺栓预紧力是静态、一次设定的；多机内力是动态、可实时调节的（接近时小、抬起时大、放下时归零），且必须保持在摩擦锥内否则打滑——螺栓没有“打滑”这个失效模式。
\end{note}

\subsection{Moore--Penrose 解为何“挤压”——Walker 的洞察}
\label{sec:cm-internal-walker}
设计者常想要“零内力”分配——既能搬物体、又不在物体里造任何挤压。直觉上 $\mathbf{f}=\mathbf{G}^{+}\mathbf{w}_{\text{ext}}$（取 $\boldsymbol{\eta}=\mathbf{0}$）应该零内力，因为它“最小范数”。但 Walker 等\cite{walker1988internal}（ICRA 1989\pz{源稿记 1988，正确年份为 1989}）指出这\emph{未必}是物理意义上的零内力分配。问题出在“内力”的定义依赖伪逆的加权方式：Moore--Penrose 伪逆最小化的是 $\|\mathbf{f}\|_2$（接触力欧氏范数），这在数学上对应一个特定的内力分解；但物理上“非挤压”要求的是\textbf{接触界面上没有相互挤压的分量}——这取决于接触点几何与物体刚度结构，未必与“最小欧氏范数”一致。Walker 的非挤压伪逆通过引入与抓取几何相关的加权 $\mathbf{W}$，构造加权伪逆
\begin{equation}
  \mathbf{G}_{\mathbf{W}}^{+}=\mathbf{W}^{-1}\mathbf{G}^\top(\mathbf{G}\mathbf{W}^{-1}\mathbf{G}^\top)^{-1},
  \label{eq:cm-weighted-pinv}
\end{equation}
使 $\mathbf{f}=\mathbf{G}_{\mathbf{W}}^{+}\mathbf{w}_{\text{ext}}$ 落在“无挤压子空间”内。关键结论：\textbf{“最小范数”和“非挤压”是两个不同的最优性，选哪个伪逆取决于你的目标}。工程实践中多数现代方法干脆放弃追求“伪逆自带零内力”，而是显式地把内力作为独立量设计（\cref{eq:cm-decomp} 的 $\boldsymbol{\eta}$ 或 \cref{eq:cm-int-basis} 的 $\boldsymbol{\lambda}$），把内力设为期望的夹紧值而非零——因为真实搬运本来就需要非零夹紧力防滑落。

\begin{insight}[别指望“一个聪明的伪逆”自动给你理想分配]
伪逆只决定\emph{运动力}那 $6$ 维怎么取（最小范数？加权最小？），它对\emph{内力}那 $3N-6$ 维是沉默的——内力必须由你根据抓持需求显式指定。把希望寄托在伪逆上，等于把一个需要工程判断的决策（夹多紧）外包给一个数学默认值（恰好为零或恰好最小范数），这正是物体被挤爆或脱手的认知根源。
\end{insight}

\paragraph{代码实现要点。}
分解的核心是伪逆 $\mathbf{G}^{+}$ 和零空间基 $\mathbf{E}$。$\mathbf{G}\mathbf{G}^\top$ 在抓取点不退化时是 $6\times 6$ 正定小矩阵，用 Cholesky/LDLT 稳定求逆得 $\mathbf{G}^{+}=\mathbf{G}^\top(\mathbf{G}\mathbf{G}^\top)^{-1}$；零空间投影 $\mathbf{P}_N=\mathbf{I}-\mathbf{G}^{+}\mathbf{G}$；零空间基 $\mathbf{E}$ 用 SVD 取右奇异矩阵末 $3N-6$ 列（SVD 最稳健，且直接给出奇异值便于检测退化）。合成接触力即 $\mathbf{f}=\mathbf{G}^{+}\mathbf{w}_{\text{ext}}+\mathbf{E}\boldsymbol{\lambda}$。三处常见错误：（1）误用左伪逆 $(\mathbf{G}^\top\mathbf{G})^{-1}\mathbf{G}^\top$——$\mathbf{G}^\top\mathbf{G}$ 是 $3N\times 3N$ 但秩仅 $6$，奇异，求逆给 NaN；行满秩宽矩阵必须用右伪逆。（2）把内力意图塞进 $\mathbf{w}_{\text{ext}}$——内力合旋量本应为 $0$，塞进 $\mathbf{w}_{\text{ext}}$ 等于要求物体产生不存在的运动，污染运动力；内力只能经 $\ker\mathbf{G}$ 注入。（3）用 $\mathbf{P}_N\times$（任意 $\boldsymbol{\eta}$）当内力而不归一化，导致内力大小不可控；用标准正交基 $\mathbf{E}$、$\boldsymbol{\lambda}$ 的每分量直接是该模式强度（牛），便于设定与限幅。

\begin{insight}[内力调节在真机里就是“夹紧力控制”]
正因内力位于 $\ker\mathbf{G}$、与运动力正交，工程上才能给多机搬运装一个独立的“夹紧力旋钮”：抓取阶段把内力（径向夹紧模式 $\boldsymbol{\lambda}$）从 $0$ 缓升到防滑所需值，搬运中保持，放下时降回 $0$——全程不影响物体跟踪轨迹。decPLM 的 “pinch”（捏紧）、绳索搬运的张力维持，本质都是在 $\ker\mathbf{G}$ 里调 $\boldsymbol{\lambda}$。这就是抽象“零空间”在真机上的物理化身。
\end{insight}

\begin{pitfall}[编程陷阱：误用 $\mathbf{G}^\top\mathbf{G}$ 求伪逆导致奇异]
套用“左伪逆” $(\mathbf{G}^\top\mathbf{G})^{-1}\mathbf{G}^\top$ 处理行满秩的宽 $\mathbf{G}$：$\mathbf{G}^\top\mathbf{G}$ 是 $3N\times 3N$ 但秩只有 $6$，严重奇异，求逆返回 NaN 或天文数字，整个力分配崩溃。根因：左伪逆要求列满秩；抓取矩阵 $\mathbf{G}$ 是 $6\times 3N$ 行满秩（列亏秩），必须用右伪逆 $\mathbf{G}^\top(\mathbf{G}\mathbf{G}^\top)^{-1}$（$\mathbf{G}\mathbf{G}^\top$ 才是 $6\times 6$ 可逆）。自检：打印 $\mathbf{G}\mathbf{G}^\top$ 的最小奇异值，远大于 $0$ 说明抓取点不退化、可安全求逆。
\end{pitfall}

\begin{pitfall}[概念误区：以为伪逆解 $\mathbf{G}^{+}\mathbf{w}$ 自动零内力]
“$\mathbf{G}^{+}$ 给最小范数解，范数最小当然没多余内力”——错。最小范数（数学）$\ne$ 非挤压（物理）。$\mathbf{G}^{+}\mathbf{w}$ 在 $\ker\mathbf{G}$ 上的分量确实为 $0$（它在 $\rowsp\mathbf{G}$ 里），所以按 Moore--Penrose 定义其“内力分量”是 $0$——但这个“零内力”是相对欧氏范数的，物理接触界面上是否真无挤压取决于加权方式（Walker）。更重要的是：真实搬运需要非零夹紧力防滑，你本就不该追求零内力！若盲信“伪逆 $=$ 零内力”，既可能误判系统无挤压（实则有），又会漏掉必须主动施加的防滑夹紧力，导致物体滑落。内力该取多少由摩擦锥定下界、由物体强度定上界，是 \cref{sec:cm-qp} 的优化问题，不是伪逆的副产品。
\end{pitfall}

\begin{pitfall}[思维陷阱：把内力当“误差”而非“需要主动设计的量”]
“内力不产生运动、还可能挤坏物体，越小越好、最好为零”——错。内力是抓持的必需品：零内力意味着零夹紧力，除非物体被双边刚性焊住，否则点接触/摩擦抓持下零夹紧 $=$ 立刻滑落。正确目标不是“消除内力”，而是“把内力调到恰当值”：下界 $=$ 防滑所需（摩擦锥），上界 $=$ 不损坏物体。把内力看成和运动力平起平坐的设计变量——运动力服务“物体怎么动”，内力服务“物体怎么被稳稳抓住”，两者都要主动给。自检：问“我的搬运若内力设为 $0$ 会怎样？”若答“物体滑落”说明内力是刚需；若答“无影响”（刚性焊接）才可设 $0$。
\end{pitfall}

\begin{practice}
\begin{enumerate}
  \item （$\star\star$）证明零空间投影 $\mathbf{P}_N=\mathbf{I}-\mathbf{G}^{+}\mathbf{G}$（$\mathbf{G}^{+}$ 为 Moore--Penrose 伪逆）满足幂等性 $\mathbf{P}_N^2=\mathbf{P}_N$ 与对称性 $\mathbf{P}_N^\top=\mathbf{P}_N$。再证对任意 $\mathbf{w}_{\text{ext}}$ 与任意 $\boldsymbol{\eta}$，运动力 $\mathbf{G}^{+}\mathbf{w}_{\text{ext}}$ 与内力 $\mathbf{P}_N\boldsymbol{\eta}$ 正交。（提示：用伪逆四条件 $\mathbf{G}\mathbf{G}^{+}\mathbf{G}=\mathbf{G}$、$\mathbf{G}^{+}\mathbf{G}\mathbf{G}^{+}=\mathbf{G}^{+}$、$(\mathbf{G}\mathbf{G}^{+})^\top=\mathbf{G}\mathbf{G}^{+}$、$(\mathbf{G}^{+}\mathbf{G})^\top=\mathbf{G}^{+}\mathbf{G}$。）
  \item （$\star\star\star$）实现一个内力调节器：输入期望外力旋量 $\mathbf{w}_{\text{ext}}$、期望夹紧强度（标量 $s>0$）、各接触点的内向单位方向，输出接触力 $\mathbf{f}$，使得（a）$\mathbf{G}\mathbf{f}=\mathbf{w}_{\text{ext}}$；（b）内力对应“所有接触点以强度 $s$ 径向内夹”。验证内力分量确实在 $\ker\mathbf{G}$ 里（$\mathbf{G}\mathbf{f}_{\text{int}}\approx\mathbf{0}$），并扫描 $s$ 从 $0$ 到 $50$ 牛画“夹紧强度 vs 接触力范数”曲线。
  \item （$\star\star\star$，跨章综合）一处用伪逆把多项式轨迹的“固定/自由导数”分解（minimum-snap 轨迹生成），本节用伪逆把接触力的“运动/内力”分解。用一张表对比两处伪逆的：（a）被分解的物理量；（b）行空间分量的物理意义；（c）零空间分量的物理意义；（d）谁是“任务约束”、谁是“自由设计量”。并论证：为什么这两个看似无关的问题能用同一套“伪逆 $+$ 零空间投影”数学统一——它们共同的抽象结构是什么？
\end{enumerate}
\end{practice}

% ════════════════════════════════════════════════════════════════
\section{leader--follower 力协调}
\label{sec:cm-leader}

\paragraph{动机：搬一个长沙发上楼。}
\cref{sec:cm-grasp,sec:cm-internal} 给了“力怎么分”的数学，但没回答“谁来决策”。当 $N$ 台机器人搬一个物体，总得有人知道“物体该去哪”。想象你和朋友搬长沙发上楼：你们不可能每挪一步都商量“向左转 $15$ 度、再向前 $0.3$ 米”——太慢太累。现实中通常一人（在前、看得见楼梯）当\textbf{主导}，另一人（在后）当\textbf{跟随}：主导者决定往哪走、怎么转，跟随者不需知道完整计划，只需\emph{感受沙发传来的力和运动}顺着使劲。这套“一人主导、其余跟随、靠物体传递意图”的模式就是 \textbf{leader--follower（主从）架构}。它的精妙在于：跟随者不需任何目标信息、也不需和主导者通信——\emph{沙发本身就是通信信道}。放到机器人上：leader 知道物体期望轨迹 $\mathbf{x}_d(t)$，用阻抗控制驱动物体跟踪；follower 不知 $\mathbf{x}_d(t)$，只测量物体（或自己接触点）的运动，据此推断 leader 意图、产生配合的力。

\paragraph{反面：全部机器人都当 leader 会怎样。}
一个自然反问：为什么不让每台都知道目标轨迹、都当 leader？这就是“全集中/全主导”方案，它在三处出问题。\textbf{其一，需给每台广播精确、时间同步的目标轨迹}：所有机器人必须在每时刻对“物体现在该在哪”达成比特级一致，任何延迟或不同步都会让不同机器人跟踪略异的目标，在刚性物体上立刻产生内力（\cref{sec:cm-internal} 的挤压）。\textbf{其二，需每台精确知道自己接触点相对物体的位置}：要跟踪物体轨迹就得知道“物体在这个位姿时我的接触点该在哪”，需要精确抓取几何标定。\textbf{其三，丧失“力放大”的好处}：leader 只需用很小的力“提议”一个运动方向，follower 们会对齐并放大这个力，使物体获得远大于 leader 单独所能提供的驱动力——全主导方案没有这种放大。

\begin{insight}[leader--follower：物体作为信息聚合器]
leader--follower 不是“偷懒少给 follower 信息”，而是用\emph{物体作为信息聚合器}，把“协调”从“通信问题”转化为“物理感知问题”。意图不通过网络传，而通过物体的运动传——这让系统对通信故障免疫、对机器人数量可扩展。代价是 follower 的响应有物理延迟（要先感受到运动才反应），且要求物体足够刚以快速传递意图。
\end{insight}

\paragraph{历史：从遥操作主从到协作搬运。}
leader--follower 的思想源头是\textbf{遥操作（teleoperation）}——人操作 master 设备、远端 slave 机器人跟随，二者通过通信链路或物理介质耦合（\cref{sec:cm-passivity} 会深入遥操作的无源性）。1990 年代研究者把主从思想从“人–机器人”推广到“机器人–机器人”协作搬运：Khatib 的多臂操作\cite{khatib1988object}、Caccavale 的双臂协调\cite{caccavale2008six} 里都有 leader/follower 角色雏形。近年的隐式通信搬运\cite{wang2018collaborative} 与 Carey \& Werfel 2021 的集体运输\cite{carey2021collective} 把这套思想推到极致：\emph{follower 完全不通信、不知道目标，仅靠本地测量物体运动来配合 leader}。
% 原句把"Force-ANTS / Force-Amplifying N-robot Transport System / leader 小力被 follower 放大"归于 Carey \& Werfel——经核为张冠李戴，已注释并改写。
\rebuilt{两处归属待修：(1)"Force-ANTS / Force-Amplifying N-robot Transport System"及"leader 小力被 follower 对齐放大"实为 Wang \& Schwager 2016 IJRR 35(13):1564--1586 的成果，\emph{非}Carey \& Werfel 2021，亦非 wang2018collaborative 所指文；(2)wang2018collaborative 题/venue(Frontiers 2018)对应论文实为 Bechlioulis \& Kyriakopoulos，作者归属不自洽。Carey \& Werfel 2021 ICRA 的真实机制是\emph{自适应柔顺}下的隐式协调，而非力放大。}这条线证明 leader--follower 可以做到极简——follower 只需一个力传感器和一条阻抗/导纳律。

\subsection{leader：用阻抗领航}
\label{sec:cm-leader-law}
leader 知道物体期望轨迹 $\mathbf{x}_d(t)$。它对物体施加阻抗律，让物体表现得像被一根弹簧–阻尼器拉向期望轨迹：
\begin{equation}
  \mathbf{w}_{\text{leader}}=\mathbf{K}_L(\mathbf{x}_d-\mathbf{x}_o)+\mathbf{D}_L(\dot{\mathbf{x}}_d-\dot{\mathbf{x}}_o),
  \label{eq:cm-leader}
\end{equation}
其中 $\mathbf{x}_o,\dot{\mathbf{x}}_o$ 是物体（leader 可观测部分）的实际位姿与速度，$\mathbf{K}_L,\mathbf{D}_L$ 是 leader 的刚度/阻尼。这本质是单臂阻抗律施加在物体上——leader 把物体“拉”向期望轨迹，但用柔性的力而非硬位置命令（遇障碍/失配时不会硬顶）。leader 通过其接触点把 $\mathbf{w}_{\text{leader}}$ 注入物体（再经 \cref{sec:cm-internal} 分解出该 leader 的接触力）。

\subsection{follower：估计意图并配合}
\label{sec:cm-follower-law}
follower 不知 $\mathbf{x}_d$，只能测物体的运动（或自己接触点的力/运动）。其目标是\emph{顺着 leader 想要的方向出力、帮物体运动，同时维持抓持内力}。常见两类策略。

\textbf{策略 A：导纳型 follower（顺从物体运动）。} follower 测物体施加在其接触点的力 $\mathbf{f}_{\text{sensed}}$（含 leader 经物体传来的“拖拽”），用导纳律产生顺从运动：
\begin{equation}
  \mathbf{M}_F\ddot{\mathbf{x}}_F+\mathbf{D}_F\dot{\mathbf{x}}_F=\mathbf{f}_{\text{sensed}}-\mathbf{f}_{\text{int,desired}},
  \label{eq:cm-follower-adm}
\end{equation}
即 follower 像一个“被推着走”的质量–阻尼系统：leader 拖拽物体、物体拖拽 follower，follower 顺从地动起来，从而不阻碍 leader 的运动，同时减去期望内力部分以维持夹紧。

\textbf{策略 B：力放大 follower（对齐并放大 leader 的力）。} 这是 Force-ANTS 的核心\rebuilt{Force-ANTS 力放大机制应归 Wang \& Schwager 2016 IJRR 35(13):1564--1586，原稿误引为 carey2021collective(Carey \& Werfel 2021，机制实为自适应柔顺)；待把此处引用改为 Wang \& Schwager 的正确条目}。follower 估计物体的运动方向 $\hat{\mathbf{d}}=\dot{\mathbf{x}}_o/\|\dot{\mathbf{x}}_o\|$（物体在往哪动，就是 leader 想要的方向），然后\emph{主动沿该方向出力}：
\begin{equation}
  \mathbf{f}_{\text{follower}}=\kappa\,\hat{\mathbf{d}}+\mathbf{f}_{\text{int}},\qquad \kappa>0.
  \label{eq:cm-follower-amp}
\end{equation}
效果：leader 用一个小力起头，让物体微微朝目标方向动；每个 follower 检测到这个运动方向、沿同方向加一份力 $\kappa\hat{\mathbf{d}}$；$N-1$ 个 follower 的力全部对齐叠加，物体获得的总驱动力是 leader 的许多倍。\textbf{leader 像乐队指挥——它自己不发声（出力小），但让所有乐手（follower）的声音对齐放大。}

\paragraph{意图估计：follower 怎么“读”出 leader 想法。}
本质是\emph{从物体运动反推期望运动方向}。最简形式：物体当前速度方向就是期望方向（策略 B 的 $\hat{\mathbf{d}}$）。更精细的可用观测器估计物体加速度、或对力信号滤波去噪。已有工作\cite{wang2018collaborative} 用物体在 follower 接触点处的局部运动作为隐式信息，证明即使 follower 只测自己接触点（不需全局物体状态），也能正确配合\rebuilt{此技术论断属实，但 wang2018collaborative 键所记题/venue(Frontiers 2018)对应论文实为 Bechlioulis \& Kyriakopoulos(frobt.2018.00090)、非 Wang \& Schwager；引用键作者归属待修正}。

\begin{note}[类比的边界：follower 估计意图 vs 交谊舞引带]
follower 估计 leader 意图，像跳交谊舞时女伴跟随男伴的引带（lead）。\emph{像的部分}：跟随者不预知舞步，仅靠搭手处传来的力/运动方向实时判断并配合，主导者用很轻的引带就能带动。\emph{不像的部分}：舞伴是智能体能预判和“补偿”，机器人 follower 只做即时反应、不预判；且舞蹈引带是单接触点，多机搬运是多接触点闭链，follower 的配合还必须满足内力约束（不能把舞伴“夹碎”）。
\end{note}

\begin{table}[H]\centering\small
\caption{leader--follower 与全集中架构的信息流对比（穷举式分类）}
\label{tab:cm-lf-vs-central}
\begin{tabular}{p{3.4cm}p{4.4cm}p{4.4cm}}
\toprule
维度 & leader--follower & 全集中（所有机器人跟踪轨迹） \\
\midrule
谁知道目标轨迹 & 仅 leader & 所有机器人 \\
机器人间通信 & 可零通信（靠物体） & 需广播同步轨迹 \\
follower 所需传感 & 力/运动传感（本地） & 全局物体状态 $+$ 抓取几何标定 \\
内力协调 & leader 定运动、各自维持内力 & 易因目标不同步产生内力 \\
可扩展性 & 高（加 follower 即可） & 低（同步负担随 $N$ 增长） \\
力放大 & 有（follower 对齐放大） & 无 \\
单点故障 & leader 失效则群龙无首 & 无单一关键点（但更脆于不同步） \\
失效模式 & leader 故障 / follower 误估方向 & 通信延迟 $\to$ 内力爆炸 \\
\bottomrule
\end{tabular}
\end{table}

\begin{insight}[Force-ANTS 的力放大：用一群“弱”机器人搬重物]
正因 follower 会对齐放大 leader 的力，Force-ANTS 这类系统才能用一群“弱”机器人搬动远超单机能力的重物——leader 甚至可以是一个人（用手轻推），$N-1$ 台机器人 follower 检测到推动方向后各加一份力，把人的轻推放大成足以移动重物的合力。这就是“力放大 $N$-robot 运输系统”名字的由来。\pz{Force-ANTS(Force-Amplifying N-robot Transport System)出自 Wang \& Schwager 2016 IJRR；本书前文曾误系于 Carey \& Werfel，此处概念描述正确但出处以 Wang \& Schwager 为准}
\end{insight}

\begin{pitfall}[概念误区：以为 follower 必须知道物体的绝对轨迹]
“不知道目标在哪，follower 怎么知道往哪推？”——follower 不需要绝对目标，只需物体的“当前运动方向”：物体在动这件事本身就编码了 leader 的意图，物体往哪动 leader 就想往哪去，follower 沿该方向出力即可（\cref{eq:cm-follower-amp}）。意图通过物体运动隐式传递，不通过轨迹广播。误以为要广播轨迹，你就退回了全集中方案，丢掉零通信和可扩展的全部好处，还引入同步导致的内力风险。延伸：follower 估计的是“方向”（单位向量），不是“位置”——这正是它对绝对定位误差鲁棒的原因。
\end{pitfall}

\begin{pitfall}[思维陷阱：把 leader 力大小等同于物体驱动力]
“leader 是主导，它的力决定物体运动，follower 只是辅助”——在力放大架构里，物体总驱动力 $\approx$ leader 力 $+\sum$ follower 力，而 follower 力之和可远大于 leader 力。leader 的作用是“定方向、起头”，不是“提供主要动力”。把 leader 力当总驱动力，会严重低估系统能搬的重量，也会误调 leader 增益。正确思维：leader 定“往哪走”（方向/意图），follower 提供“使多大劲”（放大的力）。自检：测稳态时 leader 接触力与各 follower 接触力之比，若 follower 力显著更大说明放大机制在工作。
\end{pitfall}

\begin{pitfall}[编程陷阱：follower 用未滤波的物体速度估计方向导致抖动]
直接用 $\mathbf{x}_o$ 的数值差分速度 $\hat{\mathbf{d}}=(\mathbf{x}_o[k]-\mathbf{x}_o[k-1])/\Delta t$ 归一化后用：物体近似静止或低速时，数值差分速度被噪声主导，$\hat{\mathbf{d}}$ 方向剧烈乱跳，follower 力方向随之抖动，物体震颤甚至失稳。根因：速度方向 $=$ 速度/速度模，当速度模接近 $0$ 时方向对噪声极度敏感（$0/0$ 型病态）。对策：（1）对速度低通滤波后再算方向；（2）速度模低于阈值时冻结上一方向或令 follower 力为 $0$（死区）；（3）用观测器（如卡尔曼）估计平滑速度。自检：让物体静止，观察 $\hat{\mathbf{d}}$ 是否稳定不乱跳。
\end{pitfall}

\begin{practice}
\begin{enumerate}
  \item （$\star\star$）一个 leader $+$ 三个 follower 搬运。leader 用 \cref{eq:cm-leader} 出力，follower 用 \cref{eq:cm-follower-amp} 力放大，各 $\kappa=20$ N。设稳态匀速、所有 follower 力完美对齐运动方向。若 leader 稳态出力 $10$ N，物体获得的总驱动力是多少？要把一个需 $100$ N 驱动力的重物匀速搬动，至少需几个 follower（leader 力固定 $10$ N、每 follower $\kappa=20$ N）？
  \item （$\star\star\star$）在 2D 仿真实现 leader--follower 搬运：一根刚性杆，leader 在左端用阻抗律 \cref{eq:cm-leader} 跟踪一条直线轨迹，一个 follower 在右端用力放大律 \cref{eq:cm-follower-amp}（只读杆的速度、加滤波）。验证：（a）物体跟上轨迹；（b）关闭 follower 时 leader 单独搬运更慢/更吃力，开启后更轻快（量化 leader 出力下降）。处理好速度方向抖动。
  \item （$\star\star\star$）leader--follower 的致命弱点是 leader 单点故障。设计一个“leader 选举 $+$ 切换”机制：当 follower 检测到物体长时间不动（leader 可能故障）时自动选举新 leader（如约定编号最小的健康机器人）。讨论：切换瞬间如何避免内力突变？这与共识算法（选举/一致性，\cref{ch:mr-consensus}）有什么联系？
\end{enumerate}
\end{practice}

% ════════════════════════════════════════════════════════════════
\section{去中心化力控：物体即共享黑板}
\label{sec:cm-decentral}

\paragraph{动机：蚂蚁搬食物。}
leader--follower 仍有一个“特殊角色”（leader）。能不能更彻底——让所有机器人\emph{完全同构、平等、互不通信}，仅靠各自感受物体来协调？想想一群蚂蚁搬一块远比单只大的食物：没有蚂蚁是“指挥官”、没有蚂蚁知道“全局计划”、它们也不交换复杂信息，但食物却能被搬向蚁巢。答案是：\emph{每只蚂蚁都遵循同一条简单规则}（大致是“朝蚁巢方向拽，同时感受食物运动、顺着它调整”）。当多数蚂蚁朝同一方向使劲时，少数方向不一致的蚂蚁会感受到食物被拽向某方向，于是顺势调整自己的用力方向——食物的运动成了所有蚂蚁之间唯一的“通信”。这种“所有个体同构、规则相同、靠被操作物体隐式协调”的模式就是\textbf{去中心化力控}。它与 leader--follower 的根本区别：后者有角色分化（一个 leader、若干 follower，控制律不同），去中心化\emph{所有 agent 跑完全相同的控制律}，没有任何特殊节点。

\paragraph{反面：坚持用显式通信会怎样。}
去中心化的对立面是“显式通信协调”——机器人之间用网络交换状态，由某种协议（或中央节点）算出协调的力，它在三处受限。\textbf{其一，通信延迟与带宽是硬瓶颈}：力控环常跑 $1$ kHz，要在 $1$ ms 内完成“采集$\to$广播$\to$接收$\to$协调计算$\to$下发”，对无线网络压力巨大；延迟超过控制周期则协调滞后于物理，重则因延迟注入能量而失稳（\cref{sec:cm-passivity}）。物体作为“通信信道”则是\emph{零延迟}的——力以声速在刚体中传播，远快于任何无线网络。\textbf{其二，可扩展性受限}：显式通信的协调计算常随 $N$ 增长（广播 $O(N)$、集中求解 $O(N^3)$），去中心化的本地控制律对每个 agent 是 $O(1)$ 的。\textbf{其三，单点与通信故障}：中央协调器是单点故障、无线通信可能丢包被干扰；去中心化没有中央节点、不依赖通信链路，单台故障只是少一份力，系统降级而不崩溃。

\begin{insight}[协调“涌现”于物理耦合]
去中心化力控把“协调”这个看似需要全局信息的问题，分解成 $N$ 个\emph{只用本地信息}的相同子问题，而全局协调“涌现（emerge）”于所有 agent 通过物体的物理耦合。这与共识算法的精神一致——共识让每个节点只和邻居通信就达成全局一致；去中心化搬运更进一步，连“和邻居通信”都省了，\emph{物体就是所有 agent 共同的邻居}。协调不是被某个节点计算出来的，而是从物理耦合中长出来的。
\end{insight}

\paragraph{历史：从行为群体到隐式通信。}
去中心化多机搬运根植于\textbf{群体机器人（swarm robotics）}与\textbf{基于行为的控制}——1990 年代研究的“简单规则涌现复杂群体行为”。Carey--Werfel 2021\cite{carey2021collective} 与隐式通信工作\cite{wang2018collaborative} 把它精确化为\textbf{隐式通信（implicit communication）}框架：明确指出物体的运动是 agent 间唯一的信息载体，并给出收敛性分析\rebuilt{此处"Force-ANTS"标签误系于 Carey--Werfel——Force-ANTS 属 Wang \& Schwager 2016 IJRR；且 wang2018collaborative 键的 Frontiers 2018 文实为 Bechlioulis \& Kyriakopoulos（见前注）}。Culbertson--Schwager 2020\cite{culbertson2021decentralized} 的去中心化自适应控制证明即使\emph{不知道物体质量与惯性}，每个 agent 用本地自适应律也能协同搬运（继承 Verginis 2019\cite{verginis2019cooperative} 的“无需负载参数”思想）。最新的 decPLM（Pandit 2025）\cite{pandit2025decplm} 用强化学习在 IsaacLab 里训练出零通信的 pinch-lift-move 策略，把 2--10 台 Go2 的去中心化搬运做到了实物。

\subsection{同构本地控制律}
\label{sec:cm-decentral-law}
每个 agent $i$ 运行\emph{完全相同}的控制律，只用本地量。一个有代表性的结构：
\begin{equation}
  \mathbf{f}_i=\underbrace{\mathbf{f}_i^{\text{nav}}}_{\text{导航分量}}
  +\underbrace{\mathbf{f}_i^{\text{comply}}}_{\text{顺从分量}}
  +\underbrace{\mathbf{f}_i^{\text{int}}}_{\text{内力分量}}.
  \label{eq:cm-decentral}
\end{equation}
\begin{itemize}
  \item \textbf{导航分量} $\mathbf{f}_i^{\text{nav}}$：每个 agent 自己知道（或感知）大致目标方向（如都朝目标点），独立产生一份朝向目标的力。这是“共同目标”的去中心化体现——不需协调，因为所有 agent 的目标方向大体一致。
  \item \textbf{顺从分量} $\mathbf{f}_i^{\text{comply}}$：agent 感受物体的实际运动（被其他 agent 影响的结果），用导纳律顺从它，避免与群体“对抗”。这是隐式协调的核心——它让少数方向不一致的 agent 自动向多数妥协。
  \item \textbf{内力分量} $\mathbf{f}_i^{\text{int}}$：维持抓持/夹紧，由本地约定的内力意图产生（\cref{sec:cm-internal}）。
\end{itemize}

\subsection{隐式通信：物体如何传递信息}
\label{sec:cm-decentral-implicit}
关键机制：agent $i$ 出的力 $\to$ 改变物体运动 $\to$ agent $j$ 在自己接触点测到这个运动变化 $\to$ agent $j$ 据此调整自己的力。物体把“agent $i$ 的动作”翻译成“物体运动”，再翻译成“agent $j$ 的观测”。信息流是：
\begin{equation}
  \mathbf{f}_i\;\xrightarrow{\;\mathbf{G}\;}\;\mathbf{w}_o\;\xrightarrow{\text{物体动力学}}\;\mathbf{v}_o\;\xrightarrow{\;\mathbf{G}_j^\top\;}\;\dot{\mathbf{x}}_j\;\xrightarrow{\text{agent }j\text{ 测量}}\;\mathbf{f}_j.
  \label{eq:cm-implicit-flow}
\end{equation}
注意这条链里没有任何网络通信——全是物理量经抓取矩阵的正向（$\mathbf{G}$）与对偶（$\mathbf{G}_j^\top$）映射。\textbf{物体扮演了一个零延迟、无限带宽、永不丢包的广播信道}，代价是它只能传递“物体的运动”这一种信息（不能传任意比特）。

\paragraph{收敛性的直觉。}
为什么这套各自为政的规则能收敛到协调搬运而不是一团乱？顺从分量 $\mathbf{f}_i^{\text{comply}}$ 充当\emph{负反馈}——任何与群体不一致的力都会引起物体运动，该运动被 agent 自己测到后，导纳律驱使它顺从、减少对抗。这类似共识的“误差驱动收敛”：共识里节点向邻居均值靠拢，这里 agent 向物体运动（群体合意的物理体现）靠拢。严格收敛性需要无源性论证（\cref{sec:cm-passivity}）——每个 agent 的导纳律无源、物体无源，反馈互联保持无源，从而稳定收敛。

\begin{note}[类比的边界：去中心化搬运 vs 一群人推抛锚汽车]
去中心化搬运的物体，像一群人合力推抛锚汽车时的“车”。\emph{像的部分}：没人指挥，每人朝大致前方推，靠车的运动感知整体节奏并调整自己的力，多数人方向一致时少数人自然被带动。\emph{不像的部分}：人能看到彼此、能喊话（有额外通信通道），机器人去中心化严格只靠车（物体）这一个通道；且人会预判，机器人只即时反应。把“人推车”的协调能力高估到机器人上，会让你低估去中心化对“物体刚度、感知质量”的苛刻要求。
\end{note}

\begin{table}[H]\centering\small
\caption{三种“谁来决策”架构总对比（贯穿本章的核心分类）}
\label{tab:cm-three-arch}
\begin{tabular}{p{2.6cm}p{3.2cm}p{3.2cm}p{3.2cm}}
\toprule
维度 & 全集中 & leader--follower & 去中心化 \\
\midrule
角色 & 中央协调器 $+$ 同构执行 & 1 leader $+$ $N{-}1$ follower & 全同构 \\
谁知道目标 & 中央 & leader & 所有（仅大致方向） \\
通信 & 重（状态$\leftrightarrow$中央） & 轻或零（靠物体） & 零（靠物体） \\
协调计算 & $O(N^3)$ 集中 & leader $O(1)$ & 各 $O(1)$ \\
可扩展性 & 差 & 中 & 优 \\
鲁棒性（单点） & 差（中央） & 中（leader） & 优（无关键点） \\
精度/最优性 & 高（全局信息） & 中 & 较低（仅本地信息） \\
代表工作 & 经典操作空间控制 & Force-ANTS, Wang--Schwager & decPLM, Culbertson--Schwager \\
适用 & 少机、高精度、可靠通信 & 人引导、半自主 & 大规模、零通信、强鲁棒 \\
\bottomrule
\end{tabular}
\end{table}

\begin{insight}[架构是权衡谱，不是优劣排序]
这三种架构不是“谁更先进”，而是\emph{信息可得性与系统规模的权衡谱}。全局信息可得且 $N$ 小 $\to$ 集中式拿最优；信息受限或 $N$ 大 $\to$ 去中心化拿可扩展与鲁棒，代价是放弃部分最优性。这与多智能体路径规划（MAPF）的“耦合（最优但不可扩展）vs 解耦（可扩展但次优）”是同一条权衡轴在力控层的回响。
\end{insight}

\begin{insight}[decPLM 的零通信为何可行]
正因物体是零延迟的隐式信道，decPLM\cite{pandit2025decplm} 才能让 2--10 台 Go2 在\emph{完全不通信}下协同搬运——每台 Go2 跑同一个 RL 策略，输入只有本体感知（自身姿态、足端力/接触），通过物体运动隐式感知其他 Go2 的努力。“pinch-lift-move”三阶段（捏紧$\to$抬起$\to$移动）对应 \cref{eq:cm-decentral} 里内力分量（pinch）和导航分量（move）的时序组合\pz{此为用经典力分解框架对 decPLM 的\emph{解读性}映射，非原文主张：decPLM 原文的三阶段是按命令池逐步扩张的\emph{训练课程}(pinch=接触建立、lift=加高度命令、move=加平面速度)，并用"星座对齐"奖励，全文不含抓取矩阵/$\ker\mathbf{G}$/内力旋量等表述}。零通信不是魔法，而是把通信负担转移给了物理。
\end{insight}

\begin{pitfall}[概念误区：以为“去中心化”等于“无协调/各干各的”]
“不通信、各跑各的控制律，那不就是 $N$ 个独立机器人？”——去中心化 $\ne$ 无耦合。机器人通过物体强耦合，这正是协调的来源：它们不通过“网络”协调，但通过“物体”强烈相互影响；顺从分量让每个 agent 对群体合力做负反馈，协调从物理耦合中涌现。“独立”的只是控制律的计算，不是物理上的相互作用。误以为各干各的，你会漏掉顺从分量（隐式协调的核心），写出真正独立的控制律，结果机器人互相对抗、物体被撕扯。延伸：去中心化系统的全局行为是“涌现”的，不是任何单个 agent “计算”出来的。
\end{pitfall}

\begin{pitfall}[思维陷阱：认为去中心化总优于集中式]
“零通信、可扩展、鲁棒全是优点，去中心化碾压集中式”——去中心化用“放弃全局信息”换“可扩展与鲁棒”，代价是精度与最优性下降：仅靠本地信息无法做出全局最优的力分配。少机、高精度、有可靠通信的场景（如工厂里两台标定好的机械臂精密装配），集中式用全局信息能做得更好。架构选择是权衡，不是优劣排序。按“信息可得性 $\times$ 规模 $\times$ 精度需求”选：信息全 $+$ 小规模 $+$ 高精度 $\to$ 集中；信息受限或大规模或要强鲁棒 $\to$ 去中心化。
\end{pitfall}

\begin{pitfall}[编程陷阱：所有 agent 跑同一控制律却用了不一致的坐标系]
“同构控制律”实现时各 agent 在自己机体系下计算导航力/顺从力，但没把它们统一到物体系或世界系，误以为“代码相同就是同构”。现象：各 agent 的“朝目标”方向在不同坐标系下指向不同的物理方向，合力相互抵消或乱指，物体原地打转或乱走，看起来像没协调。根因：“同构”指物理行为同构（都朝同一物理目标使劲），不是“代码字节相同”——坐标系不对齐时相同代码产生不同物理方向。对策：明确所有力的参考系，导航方向应在共同参考系（世界系/物体系）中定义，再变换到各 agent 执行系。自检：让所有 agent 目标设为同一世界方向，检查各接触力在世界系下是否真同向。
\end{pitfall}

\begin{practice}
\begin{enumerate}
  \item （$\star\star$）解释为什么去中心化搬运中“顺从分量” $\mathbf{f}_i^{\text{comply}}$ 是协调收敛的关键。若去掉顺从分量、只保留导航分量和内力分量会发生什么？用“两台机器人导航方向有 $30$ 度夹角”（一台想往北、一台想往东北）的具体例子说明。
  \item （$\star\star\star$）在 2D 仿真实现 $N=3$ 去中心化搬运：每台跑 \cref{eq:cm-decentral} 的同一控制律（导航分量朝同一目标点 $+$ 顺从分量对物体速度做导纳 $+$ 固定内力），三台初始导航方向因定位噪声略有不同。验证：（a）物体最终被搬向目标；（b）关闭顺从分量后三台“对抗”、物体偏离或抖动。务必处理坐标系对齐。
  \item （$\star\star\star$，跨章综合）论述“去中心化”思想在多机系统三处的体现：共识（节点靠邻居通信达成一致，\cref{ch:mr-consensus}）、去中心化 MAPF/PIBT（机器人靠局部规则避碰）、本章去中心化搬运（agent 靠物体隐式协调）。三者的“局部信息”“协调机制”“涌现的全局行为”分别是什么？它们共享的权衡（放弃什么换取什么）是什么？
\end{enumerate}
\end{practice}

% ════════════════════════════════════════════════════════════════
\section{协调阻抗与物体阻抗}
\label{sec:cm-impedance}

\paragraph{动机：物体撞到障碍时，谁该“软”下来。}
前面解决了“力怎么分”（\cref{sec:cm-internal}）和“谁决策”（\cref{sec:cm-leader,sec:cm-decentral}），但还没说清“用什么力学行为搬”。单臂力控里我们学过阻抗控制（让末端表现得像弹簧–阻尼器）。多机搬运里阻抗该加在哪？两台机器人搬一块平板，平板前缘撞上墙，我们希望系统柔顺——不要硬顶（会损坏机器人、墙或物体），而要像弹簧一样“知难而退”。问题是：阻抗（柔顺）应体现在哪个层面？\emph{选项一：各末端独立阻抗}——每台末端各跑一套阻抗律，各自对自己接触点的位置偏差产生回复力；问题是两套独立阻抗各管各的末端，对“物体整体该怎么柔顺响应碰撞”没有共识，平板撞墙时靠墙侧机器人感到大力、另一侧没感觉，响应不协调还易产生内力。\emph{选项二：物体阻抗}——把阻抗加在\textbf{物体质心}上，让物体整体表现得像一个有期望惯量 $\mathbf{M}_d$、阻尼 $\mathbf{D}_d$、刚度 $\mathbf{K}_d$ 的刚体：物体受外力（撞墙）时按此期望阻抗“退让”；控制器先算出“为实现物体的期望阻抗行为，物体需承受多大合力旋量 $\mathbf{w}_{\text{ext}}$”，再用 \cref{sec:cm-internal} 的分解把 $\mathbf{w}_{\text{ext}}$ 分给各机器人。我们关心的是物体的行为，就该把期望行为定义在物体上——这就是 object impedance 的核心思想。

\paragraph{反面：用各末端独立阻抗会怎样。}
坚持选项一会暴露三个问题。\textbf{其一，碰撞响应不协调}：外部扰动作用在物体某处，但只有附近末端“感觉到”并柔顺退让，远处末端不退让，物体响应取决于碰撞位置、不可预测；物体阻抗则让整个物体按统一 $\mathbf{M}_d,\mathbf{D}_d,\mathbf{K}_d$ 响应，与碰撞位置无关。\textbf{其二，内力与运动阻抗纠缠}：各末端阻抗回复力里既含“让物体运动”的分量、也含“末端间相互挤压”的分量，二者揉在一起，调阻抗参数时想让物体更柔顺却意外改变了内力；物体阻抗 $+$ \cref{sec:cm-internal} 分解则把运动阻抗（物体级 $\mathbf{K}_d,\mathbf{D}_d$）与内力（$\ker\mathbf{G}$ 里独立设定）彻底解耦。\textbf{其三，期望位置不一致即内力}：各末端独立阻抗需给每个末端一个期望位置 $\mathbf{x}_{d,i}$，这些必须严格满足刚体约束（都来自同一物体期望位姿），否则不一致的期望位置在刚性物体上直接产生内力；物体阻抗只需一个物体期望位姿 $\mathbf{x}_{o,d}$，各末端位置由刚体几何自动导出、天然一致。

\begin{insight}[在“任务真正所在的空间”定义期望行为]
物体阻抗与各末端独立阻抗的区别，本质是\emph{“在哪个空间定义期望行为”}。各末端阻抗在 $3N$ 维接触空间定义期望（过定义，必然引入内力纠缠）；物体阻抗在 $6$ 维物体空间定义期望（恰好定义，运动与内力干净分离）。永远在“任务真正所在的空间”（这里是 $6$ 维物体空间）定义期望行为，再映射到执行器空间——这是力控的一条普适原则，与操作空间控制“在任务空间定义行为”一脉相承。
\end{insight}

\paragraph{历史：从对称构型到物体级阻抗。}
Caccavale、Chiacchio 等在 1990 年代为双臂协调提出\textbf{对称构型公式（symmetric formulation）}：用“绝对运动”（物体质心位姿）和“相对运动”（两末端相对位姿，编码内力）两组变量描述双臂，把阻抗加在绝对运动上、力调节加在相对运动上\cite{caccavale2008six}。这正是物体阻抗的早期形态——绝对运动 $=$ 物体运动、相对运动 $=$ 内力维度。Wimb\"ock、Ott、Hirzinger 2012\cite{wimbock2012comparison} 在 DLR 系统比较了多种物体级抓取控制器，提出统一的物体级阻抗框架，把多指手和多臂统一处理，并明确了内力沿抓取线（internal force along grasp lines）的设计。这条线确立了“物体阻抗 $+$ 内力调节”作为多机/多指柔顺操作的标准范式。

\subsection{期望物体阻抗行为}
\label{sec:cm-impedance-desired}
我们希望被搬物体表现得像一个目标阻抗系统。设物体期望位姿 $\mathbf{x}_{o,d}$、实际位姿 $\mathbf{x}_o$，跟踪误差 $\tilde{\mathbf{x}}_o=\mathbf{x}_o-\mathbf{x}_{o,d}$。期望阻抗动力学（在物体质心、$6$ 维）：
\begin{equation}
  \mathbf{M}_d\ddot{\tilde{\mathbf{x}}}_o+\mathbf{D}_d\dot{\tilde{\mathbf{x}}}_o+\mathbf{K}_d\tilde{\mathbf{x}}_o=\mathbf{w}_{\text{ext,env}},
  \label{eq:cm-obj-imp}
\end{equation}
其中 $\mathbf{w}_{\text{ext,env}}$ 是环境施加给物体的外部力旋量（碰撞、人推等）。这条式子说：物体偏离期望轨迹时受弹簧 $\mathbf{K}_d$ 的回复、阻尼 $\mathbf{D}_d$ 的耗散；遇环境力时按期望惯量 $\mathbf{M}_d$ 的“软硬”退让。$\mathbf{M}_d,\mathbf{D}_d,\mathbf{K}_d$ 是 $6\times 6$ 矩阵，整形物体整体柔顺行为。这与 Hogan 阻抗律形式完全相同，只是作用对象从“末端”变成“物体质心”。

\subsection{求所需物体合力旋量}
\label{sec:cm-impedance-wrench}
物体的真实动力学（刚体）为
\begin{equation}
  \mathbf{M}_o\ddot{\mathbf{x}}_o+\mathbf{C}_o(\dot{\mathbf{x}}_o)=\mathbf{w}_o^{\text{robots}}+\mathbf{w}_{\text{ext,env}}+\mathbf{w}_{\text{grav}},
  \label{eq:cm-obj-dyn}
\end{equation}
其中 $\mathbf{M}_o$ 是物体的 $6\times 6$ 空间惯量矩阵，$\mathbf{C}_o$ 含科氏/离心项，$\mathbf{w}_o^{\text{robots}}=\mathbf{G}\mathbf{f}$ 是所有机器人施加的合力旋量，$\mathbf{w}_{\text{grav}}$ 是重力旋量。把 \cref{eq:cm-obj-imp} 解出期望加速度
\[
  \ddot{\mathbf{x}}_o^{\text{des}}=\ddot{\mathbf{x}}_{o,d}-\mathbf{M}_d^{-1}\big(\mathbf{D}_d\dot{\tilde{\mathbf{x}}}_o+\mathbf{K}_d\tilde{\mathbf{x}}_o-\mathbf{w}_{\text{ext,env}}\big),
\]
代入 \cref{eq:cm-obj-dyn} 求 $\mathbf{w}_o^{\text{robots}}$：
\begin{equation}
\boxed{\;
  \mathbf{w}_{\text{ext}}:=\mathbf{w}_o^{\text{robots}}=\mathbf{M}_o\ddot{\mathbf{x}}_o^{\text{des}}+\mathbf{C}_o(\dot{\mathbf{x}}_o)-\mathbf{w}_{\text{grav}}-\mathbf{w}_{\text{ext,env}}.
\;}
  \label{eq:cm-wrench-needed}
\end{equation}
这就是物体级控制器要让机器人共同施加的外力旋量。它需要物体动力学参数 $\mathbf{M}_o,\mathbf{C}_o$——这是\textbf{基于模型}的物体阻抗。若物体参数未知，可改用 Verginis/Culbertson 的\emph{自适应}方案（在线估计 $\mathbf{M}_o$）或去中心化方案，避开对精确模型的依赖。

\subsection{完整管线：物体阻抗 \texorpdfstring{$\to$}{->} 接触力}
\label{sec:cm-impedance-pipeline}
把 \cref{eq:cm-wrench-needed} 算出的 $\mathbf{w}_{\text{ext}}$ 经 \cref{sec:cm-internal} 分解、加上独立设定的内力，得到各机器人接触力：
\begin{equation}
  \mathbf{f}=\mathbf{G}^{+}\mathbf{w}_{\text{ext}}+\mathbf{E}\boldsymbol{\lambda}_{\text{int}}.
  \label{eq:cm-imp-compose}
\end{equation}
然后每台机器人 $i$ 取出自己那 $3$ 维 $\mathbf{f}_i$，通过雅可比转置 $\boldsymbol{\tau}_i=\mathbf{J}_i^\top\mathbf{f}_i$ 转成自己的关节力矩执行。完整数据流如下（每一步标注公式）。

\begin{center}
\small
\begin{tikzpicture}[
  node distance=4mm and 6mm, >=Stealth,
  blk/.style={draw, rounded corners, align=center, font=\small, inner sep=3pt, fill=main!6, draw=main!60!black},
  ln/.style={->, thick, gray!60!black}, lbl/.style={font=\scriptsize, gray!50!black}]
\node[blk] (xd) {物体期望轨迹 $\mathbf{x}_{o,d}$};
\node[blk, below=7mm of xd] (acc) {期望物体加速度 $\ddot{\mathbf{x}}_o^{\text{des}}$};
\node[blk, below=7mm of acc] (w) {所需合力旋量 $\mathbf{w}_{\text{ext}}$};
\node[blk, below=7mm of w] (f) {各接触力 $\mathbf{f}_i$};
\node[blk, below=7mm of f] (tau) {各机器人关节力矩 $\boldsymbol{\tau}_i$};
\draw[ln] (xd)--(acc) node[midway,right,lbl]{阻抗律 \cref{eq:cm-obj-imp}};
\draw[ln] (acc)--(w) node[midway,right,lbl]{物体动力学 \cref{eq:cm-wrench-needed}};
\draw[ln] (w)--(f) node[midway,right,lbl]{分解 \cref{eq:cm-imp-compose}（含内力意图 $\boldsymbol{\lambda}_{\text{int}}$）};
\draw[ln] (f)--(tau) node[midway,right,lbl]{雅可比转置 $\boldsymbol{\tau}_i=\mathbf{J}_i^\top\mathbf{f}_i$};
\end{tikzpicture}
\end{center}

\begin{note}[类比的边界：物体阻抗 vs 汽车悬架]
物体阻抗让被搬物体像“装在虚拟弹簧–阻尼悬架上的刚体”。\emph{像的部分}：物体对外部扰动的响应由 $\mathbf{M}_d,\mathbf{D}_d,\mathbf{K}_d$ 整形，像悬架决定车身对路面冲击的响应——刚度大则硬、阻尼大则不晃。\emph{不像的部分}：真实悬架是被动元件（参数固定、本征无源），物体阻抗是主动控制（参数可在线变，变阻抗会有无源性风险，见 \cref{sec:cm-passivity}）；且悬架只在一个方向，物体阻抗是 $6$ 维全姿态。把“悬架天然稳定”照搬到主动变阻抗，会忽略 \cref{sec:cm-passivity} 要解决的能量注入危险。
\end{note}

\subsection{内力阻抗：连内力也可以是柔顺的}
\label{sec:cm-impedance-internal}
更精细的框架（Wimb\"ock--Ott\cite{wimbock2012comparison}）不仅给物体运动加阻抗，也给\textbf{内力加阻抗}——内力不是设成硬值，而是让内力误差也表现出弹簧–阻尼行为：$\boldsymbol{\lambda}_{\text{int}}$ 由内力跟踪误差经一个内力阻抗律产生。好处：抓持也变柔顺，物体被外力顶时夹紧力能平滑调整而非硬抗。这把“运动柔顺”和“抓持柔顺”统一在阻抗框架下，二者分别作用在 $\range(\mathbf{G}^\top)$（运动）和 $\ker\mathbf{G}$（内力）正交子空间，互不干扰。

\paragraph{代码实现要点。}
一拍物体阻抗控制：（1）阻抗律解期望物体加速度 $\ddot{\mathbf{x}}_o^{\text{des}}=\ddot{\mathbf{x}}_{o,d}-\mathbf{M}_d^{-1}(\mathbf{D}_d\dot{\tilde{\mathbf{x}}}_o+\mathbf{K}_d\tilde{\mathbf{x}}_o-\mathbf{w}_{\text{env}})$；（2）由物体动力学求所需合力旋量 $\mathbf{w}_{\text{ext}}=\mathbf{M}_o\ddot{\mathbf{x}}_o^{\text{des}}+\mathbf{C}_o-\mathbf{w}_{\text{grav}}-\mathbf{w}_{\text{env}}$（准静态/低速可略 $\mathbf{C}_o$）；（3）复用 \cref{sec:cm-internal} 的分解，返回 $\mathbf{G}^{+}\mathbf{w}_{\text{ext}}+\mathbf{E}\boldsymbol{\lambda}_{\text{int}}$。常见错误：把环境力 $\mathbf{w}_{\text{env}}$ 符号搞反（物体撞墙时不退让反而“主动顶上去”，危险），或漏掉重力补偿（机器人没补偿重力，物体被搬起后整体下坠，控制器拼命用运动力补、内力被挤占，抓持可能滑脱）。

\begin{insight}[DLR 的多指手和双臂用同一套物体阻抗]
正因物体阻抗把期望行为定义在被操作物体上、与接触点数量解耦，DLR 才能用\emph{同一套物体阻抗框架}统一处理多指手抓取（$N$ 个指尖）和双臂搬运（$2$ 个末端）——区别只在抓取矩阵 $\mathbf{G}$ 的尺寸与内力维度，控制律结构（\crefrange{eq:cm-obj-imp}{eq:cm-imp-compose}）完全一致。你在本章学的物体阻抗，原样适用于双臂、多指灵巧手、多足搬运。
\end{insight}

\begin{pitfall}[概念误区：把物体阻抗等同于“每个末端各跑一套阻抗”]
“每台机器人末端跑笛卡尔阻抗，合起来就是多机阻抗”——$N$ 个独立末端阻抗 $\ne$ 物体阻抗。前者在 $3N$ 维接触空间过定义期望，运动与内力纠缠、碰撞响应不协调；后者在 $6$ 维物体空间定义统一阻抗、再分配，运动与内力解耦。两者行为差异巨大。用 $N$ 个独立末端阻抗，你既无法保证物体整体柔顺一致，又会被期望位置不一致引入的内力反复折磨。延伸：何时各末端阻抗也可接受？当物体不是刚性闭链（如各机器人搬独立物体），或内力维度为 $0$（双机点接触）时。
\end{pitfall}

\begin{pitfall}[编程陷阱：物体阻抗漏掉重力旋量补偿]
计算 $\mathbf{w}_{\text{ext}}$ 时忘记减去/补偿物体重力 $\mathbf{w}_{\text{grav}}$（写成 $\mathbf{w}_{\text{ext}}=\mathbf{M}_o\ddot{\mathbf{x}}-\mathbf{w}_{\text{env}}$）。现象：物体被抬起后持续下坠，控制器用越来越大的运动力去补，运动力挤占接触力预算、内力（夹紧）被压缩，物体可能滑脱；或表现为稳态位置低于期望。根因：物体阻抗的运动方程必须包含所有作用在物体上的力，重力是恒定外力，不补偿就成了持续扰动。对策：显式建模重力旋量 $\mathbf{w}_{\text{grav}}=[m_o\mathbf{g}^\top,\,(\mathbf{r}_{\text{cog}}\times m_o\mathbf{g})^\top]^\top$ 并在 $\mathbf{w}_{\text{ext}}$ 中补偿。自检：物体悬停（期望静止）时，各接触力之和的竖直分量应 $\approx$ 物体重量。
\end{pitfall}

\begin{pitfall}[思维陷阱：以为加大物体刚度 $\mathbf{K}_d$ 总能提升跟踪精度]
“跟踪误差大？调大 $\mathbf{K}_d$，弹簧硬一点就跟得紧了”——$\mathbf{K}_d$ 过大有三害：（1）牺牲柔顺——撞到障碍/人时硬顶，失去力控的安全意义；（2）放大内力——刚性物体上微小误差经大 $\mathbf{K}_d$ 变成大接触力，挤压加剧；（3）逼近稳定边界——大刚度需大阻尼配合否则振荡，且对建模误差敏感。物体阻抗的价值恰在“柔”，一味加 $\mathbf{K}_d$ 等于退回硬位置控制。正确思维：$\mathbf{K}_d$ 按任务选，精密对接需较大但仍有限的刚度 $+$ 足够阻尼，人机协作/避撞需小刚度；跟踪误差大应先查前馈（$\ddot{\mathbf{x}}_{o,d}$）、重力补偿、模型精度，而非无脑加 $\mathbf{K}_d$。
\end{pitfall}

\begin{practice}
\begin{enumerate}
  \item （$\star\star$）从期望物体阻抗 \cref{eq:cm-obj-imp} 出发，完整推导所需合力旋量 \cref{eq:cm-wrench-needed} 的每一步，明确每项（$\mathbf{M}_o\ddot{\mathbf{x}}^{\text{des}}$、$\mathbf{C}_o$、$-\mathbf{w}_{\text{grav}}$、$-\mathbf{w}_{\text{ext,env}}$）的物理意义。然后讨论：当 $\mathbf{M}_d=\mathbf{M}_o$（期望惯量取真实惯量）时式子如何简化？这种取法的物理含义（物体“感觉不到”惯量被改变）是什么？
  \item （$\star\star\star$）扩展物体阻抗一拍控制，加入“内力阻抗”——内力意图 $\boldsymbol{\lambda}_{\text{int}}$ 不再是常数，而由内力跟踪误差经一个一阶阻抗律产生（输入期望内力、测得内力，输出 $\boldsymbol{\lambda}_{\text{int}}$）。在双臂搬运里测试：给物体施加一个外部冲击，对比“硬内力”和“内力阻抗”两种方式下夹紧力的响应平滑度。
  \item （$\star\star\star$）物体阻抗 \cref{eq:cm-wrench-needed} 依赖已知物体惯量 $\mathbf{M}_o$。设计回退方案：当 $\mathbf{M}_o$ 未知时，如何用准静态假设（忽略 $\mathbf{M}_o\ddot{\mathbf{x}}$ 项，只补偿重力 $+$ 阻抗回复）退化出不需要惯量的简化物体阻抗控制器？讨论这个简化在什么速度范围内可接受、何时失效（提示：联系 Verginis 2019 “仅适用准静态”的局限）。
\end{enumerate}
\end{practice}

% ════════════════════════════════════════════════════════════════
\section{负载分配优化：摩擦锥下的 QP}
\label{sec:cm-qp}

\paragraph{动机：伪逆给出的力，机器人做不到。}
\cref{sec:cm-internal} 用伪逆把外力旋量分给各机器人，但伪逆有两个盲点：（1）它不管接触力是否物理可行——可能要求某接触点拉物体（点接触/摩擦抓持拉不了）或超出机器人力极限；（2）它最小化的是欧氏范数，未必是你真正想优化的（如让重物分给更强的机器人）。三台机器人用脚（点接触带摩擦）托举一块平板，用伪逆 $\mathbf{f}=\mathbf{G}^{+}\mathbf{w}_{\text{ext}}$ 算出分配后下发，却发现：某接触点的伪逆解要求\emph{向下拉}物体（法向力为负），但脚和平板是单边接触——脚只能往上推、不能往下拉，这个力做不到，机器人脚直接离开物体；另一个接触点的切向力远超摩擦锥，必然打滑；还有一个要求 $500$ N、超过那台机器人 $300$ N 的力极限，电机饱和、物体倾斜。伪逆对这些\textbf{物理约束一无所知}——它只解了一个无约束的最小范数问题。真实搬运里接触力必须满足：合成正确的物体旋量（等式约束）、法向力非负（单边接触）、切向力在摩擦锥内（不打滑）、不超力极限（不饱和）。这些约束让“力分配”从一个线性代数问题（伪逆）变成一个\textbf{带约束优化问题（QP）}。

\paragraph{反面：只用伪逆不加约束会怎样。}
无约束伪逆在三类情形必然出事。\textbf{其一，单边接触被违反}：点接触、脚踩、绳索牵引都是单边的（只能推/拉其中一个方向），伪逆可能给出反方向的力，接触直接断开。\textbf{其二，摩擦锥被违反、物体打滑}：摩擦能提供的切向力有上限 $\|\mathbf{f}_t\|\le\mu f_n$，伪逆的最小范数解经常把切向力用得很满（它不知道有这个限制）导致打滑；物体一滑，抓取几何 $\mathbf{r}_i$ 变了、$\mathbf{G}$ 变了，整个控制失准，常雪崩式失败。\textbf{其三，无法表达“分配偏好”}：你想让强壮的机器人多担、虚弱的少担，伪逆的“最小范数”是写死的偏好，不能表达工程意图；QP 的目标函数可自由编码这些偏好（加权）。

\begin{insight}[从“求一个解”到“在可行集里求最优解”]
从伪逆到 QP，是从“求一个解”到“在可行集里求最优解”的跃迁。伪逆假装力空间是无约束的 $\mathbb{R}^{3N}$，QP 正视力空间被摩擦锥、单边接触、力极限切割成的\emph{可行多面体/锥}。多机搬运的真正难点不在“算出一组力”，而在“算出一组\textbf{物理可实现}的力”——后者必须是约束优化。这与 MAPF 从“各自最短路”到“带冲突约束的搜索”是同一种认识升级：现实约束让简单解法失效，必须显式建模约束。
\end{insight}

\paragraph{历史：从力分配到接触力优化。}
把接触力分配写成优化问题的传统可追溯到 Orin--Oh 1986\cite{orin1986force} 的闭链力分配（已含约束意识）和抓取力优化文献（1990s--2000s 研究在摩擦锥约束下最小化抓取力）。在\textbf{腿足机器人}领域，这套方法以“接触力优化/力分配 QP”之名成为站立平衡与运动控制的标配——把躯干所需净力旋量分配到各支撑足的接触力，约束是摩擦锥（线性化为摩擦金字塔）和单边法向力。这与多机搬运的负载分配是\emph{同一个数学问题}（都是“把物体所需旋量分给多个受摩擦约束的接触点”），只是“接触点”从“足”变成“机器人末端”。现代实现普遍用 QP 求解器（qpOASES、OSQP、quadprog）在每个控制周期实时求解。

\subsection{决策变量与目标}
\label{sec:cm-qp-obj}
决策变量是所有接触力 $\mathbf{f}\in\mathbb{R}^{3N}$。一个通用目标函数：
\begin{equation}
  \min_{\mathbf{f}}\;\tfrac{1}{2}\mathbf{f}^\top\mathbf{W}\mathbf{f}+\mathbf{g}^\top\mathbf{f},
  \label{eq:cm-qp-cost}
\end{equation}
$\mathbf{W}\succ 0$ 是权重矩阵，$\mathbf{g}$ 是线性项。常见取法：$\mathbf{W}=\mathbf{I}$（最小化总接触力范数，$\approx$ 伪逆但带约束）；$\mathbf{W}=\blkdiag(w_1\mathbf{I}_3,\dots,w_N\mathbf{I}_3)$，$w_i$ 大则惩罚 agent $i$ 出力（\textbf{让弱机器人少担、强机器人多担}）；加项惩罚切向力（鼓励接近法向、远离摩擦锥边界，提升抗滑裕度）。

\subsection{等式约束：合成正确旋量}
\label{sec:cm-qp-eq}
接触力必须合成期望物体外力旋量（来自 \cref{sec:cm-impedance} 物体阻抗）：
\begin{equation}
  \mathbf{G}\mathbf{f}=\mathbf{w}_{\text{ext}}.
  \label{eq:cm-qp-eq}
\end{equation}
这把 \cref{sec:cm-internal} 的“必须满足 $\mathbf{G}\mathbf{f}=\mathbf{w}$”作为硬等式约束放进 QP。QP 会在所有满足 \cref{eq:cm-qp-eq} 的力（即 \cref{sec:cm-internal} 的解空间 $\mathbf{G}^{+}\mathbf{w}+\ker\mathbf{G}$）里挑最优的——\textbf{QP 自动地在内力子空间里优化}，因为内力正是 \cref{eq:cm-qp-eq} 的解空间自由度。

\subsection{不等式约束：摩擦锥、单边、力限}
\label{sec:cm-qp-ineq}
\textbf{单边接触（法向力非负）}：设接触点 $i$ 的法向（指向物体）为 $\mathbf{n}_i$，则法向力 $f_{i,n}=\mathbf{n}_i^\top\mathbf{f}_i\ge 0$（只能推不能拉；绳索则相反，只能拉不能推）：
\begin{equation}
  \mathbf{n}_i^\top\mathbf{f}_i\ge 0.
  \label{eq:cm-qp-unilateral}
\end{equation}
\textbf{摩擦锥（不打滑）}：切向力受库仑摩擦限制 $\|\mathbf{f}_{i,t}\|\le\mu_i f_{i,n}$。这是二阶锥约束，QP（二次规划）不能直接处理非线性锥，故\textbf{线性化为摩擦金字塔（friction pyramid）}——用 $k$ 个切向方向把圆锥近似成多面锥：
\begin{equation}
  \pm\mathbf{t}_{i,j}^\top\mathbf{f}_i\le\mu_i\,\mathbf{n}_i^\top\mathbf{f}_i,\quad j=1,\dots,k/2,
  \label{eq:cm-qp-friction}
\end{equation}
$\mathbf{t}_{i,j}$ 是切平面内的采样方向。$k$ 越大近似越精（$k=4$ 是常用方金字塔，$k=8$ 更精），代价是约束数随 $k$ 增长。\textbf{力上限（不饱和）}：$\|\mathbf{f}_i\|\le f_i^{\max}$，线性化或用框约束 $|f_{i,x}|,|f_{i,y}|,|f_{i,z}|\le f^{\max}/\sqrt{3}$ 等。把上述堆成 $\mathbf{A}\mathbf{f}\le\mathbf{b}$，得 QP 标准形：
\begin{equation}
\boxed{\;
  \min_{\mathbf{f}}\ \tfrac12\mathbf{f}^\top\mathbf{W}\mathbf{f}+\mathbf{g}^\top\mathbf{f}
  \quad\text{s.t.}\quad
  \mathbf{G}\mathbf{f}=\mathbf{w}_{\text{ext}},\;\; \mathbf{A}\mathbf{f}\le\mathbf{b}.
\;}
  \label{eq:cm-qp-std}
\end{equation}

\begin{table}[H]\centering\small
\caption{三种负载分配策略对比（穷举式分类）}
\label{tab:cm-qp-strategies}
\begin{tabular}{p{1.8cm}p{4.6cm}p{2.6cm}p{1.4cm}p{2.4cm}}
\toprule
策略 & 方法 & 约束处理 & 计算量 & 适用 \\
\midrule
等分 & $\mathbf{f}_i=\mathbf{w}_{\text{ext}}/N$（仅力，忽略力矩/几何） & 无 & 极低 & 对称布局、粗略估计 \\
加权伪逆 & $\mathbf{G}_{\mathbf{W}}^{+}=\mathbf{W}^{-1}\mathbf{G}^\top(\mathbf{G}\mathbf{W}^{-1}\mathbf{G}^\top)^{-1}$ & 无（不保证可行） & 低（闭式） & 无单边/摩擦约束、需偏好加权 \\
QP 优化 & 解 \cref{eq:cm-qp-std} & 摩擦锥 $+$ 单边 $+$ 力限 & 中（迭代求解器） & 真实接触、需保证可行性 \\
\bottomrule
\end{tabular}
\end{table}

\noindent 等分最简但常违反力矩平衡（非对称布局时物体会转）；加权伪逆能编码偏好且有闭式解，但不保证摩擦锥/单边可行；QP 最完备但需实时求解器。工程上：原型用加权伪逆快速验证，上真机用 QP 保证物理可行。

\begin{note}[类比的边界：负载分配 QP vs 几个人合抬桌子]
负载分配 QP 与“几个人合抬一张桌子，怎么站位、怎么使劲最省力又不让桌子滑”类似。\emph{像的部分}：都要满足“桌子被抬起且不转”（等式约束），每人的手不能打滑（摩擦约束）、不能用超过自己力气的劲（力限），并在可行方式里挑最舒服的（目标）。\emph{不像的部分}：人能凭经验和实时触觉自动满足这些约束，QP 必须把每条约束显式写成不等式且线性化摩擦锥；人的“省力”是模糊感觉，QP 的目标必须是精确的二次型。把人的直觉协调高估到机器人上，会让你忽略约束建模的繁琐与摩擦锥线性化的精度损失。
\end{note}

\paragraph{代码实现要点。}
用 OSQP/qpOASES 风格接口求解 \cref{eq:cm-qp-std}：目标 Hessian $\mathbf{H}=\blkdiag(w_i\mathbf{I}_3)$（弱机器人权大$\to$少出力），加 $10^{-6}\mathbf{I}$ 正则保证严格正定；等式约束 $\mathbf{A}_{\text{eq}}=\mathbf{G}$、$\mathbf{b}_{\text{eq}}=\mathbf{w}_{\text{ext}}$；不等式逐接触点构造——切向基 $\mathbf{t}_1,\mathbf{t}_2$ 与法向正交，单边 $-\mathbf{n}_i^\top\mathbf{f}\le 0$，摩擦金字塔 $\pm\mathbf{t}^\top\mathbf{f}-\mu\mathbf{n}^\top\mathbf{f}\le 0$（$4$ 边），力限框约束。三处常见错误：（1）把二阶锥 $\|\mathbf{f}_t\|\le\mu f_n$ 直接喂给 QP——QP 只接受线性不等式，二阶锥需 SOCP 求解器，必须线性化为金字塔；（2）Hessian 不正则化（纯 blkdiag 可能半正定）$\to$ QP 退化、解不唯一，加 $10^{-6}\mathbf{I}$；（3）忘记单边约束只加摩擦 $\to$ 求解器可能给出“拉”接触的解，点接触/脚/绳索物理上做不到，单边约束必不可少。

\begin{insight}[腿足 WBC 的力分配 QP 与多机搬运是同一个 QP]
正因“把物体所需旋量分给多个受摩擦约束的接触点”是同一个数学，腿足机器人 WBC（如 WBIC/OCS2）里每周期求解的足端力分配 QP，与多机搬运的负载分配 QP \cref{eq:cm-qp-std} \emph{属同一 QP 族}（同类 wrench 分配：接触力二次代价 $+$ 抓取/力旋量映射等式 $+$ 摩擦锥/单边不等式）——只是“接触点”是足、“物体”是机器人躯干。\pz{"结构完全一致"措辞偏强：WBIC 还额外耦合全身浮基动力学、且决策变量是"相对 MPC 反力的增量"，搬运 QP 则用接触幅值标量、无全身动力学；二者同源同族而非逐项等同。}掌握了本节的 QP，你既能做多机搬运的力分配，也能读懂腿足控制的接触力优化；qpOASES/OSQP 这些求解器在两个领域通用。
\end{insight}

\begin{pitfall}[编程陷阱：把二阶摩擦锥直接喂给 QP 求解器]
用 $\|\mathbf{f}_t\|\le\mu f_n$（二阶锥）作为 QP 约束：QP 求解器（只接受线性约束）报错、或静默忽略该约束给出打滑的解；用了 SOCP 求解器才能处理但算力大增。根因：二次规划（QP）的约束必须是线性的，摩擦锥 $\|\mathbf{f}_t\|\le\mu f_n$ 是二阶锥，属 SOCP 范畴。对策：线性化为摩擦金字塔（\cref{eq:cm-qp-friction}），用 $k$ 个切向方向近似圆锥，$k=4$（方金字塔）最常用、$k=8$ 更精。自检：增大 $k$ 看解是否收敛、打滑是否减少；需要精确锥时改用 SOCP 求解器。
\end{pitfall}

\begin{pitfall}[概念误区：以为加权伪逆能满足摩擦/单边约束]
“加权伪逆能编码偏好，调好权重就能避免打滑/拉力了”——加权伪逆只改变在解空间里选哪个最小（加权）范数解，它仍是无约束闭式解，不知道摩擦锥和单边约束的存在，给出的解照样可能要求拉接触或超摩擦锥。权重影响“分配偏好”，不构成“可行性保证”。误信加权伪逆可行，你会跳过 QP 直接上加权伪逆，真机一接触就打滑/脱手，还以为是权重没调好越调越乱。延伸：加权伪逆 $=$ 无约束 QP 的闭式解；一旦有不等式约束（摩擦/单边/力限），就没有闭式解，必须迭代求解 QP。
\end{pitfall}

\begin{pitfall}[思维陷阱：QP 无解时归咎于求解器而非物理不可行]
“OSQP 说无解，换个求解器/调调容差应该就好了”——QP 无解往往是物理事实：在当前接触几何、摩擦系数、力限下，根本不存在能合成所需旋量又不打滑/不超限的接触力（例如物体太重、摩擦太小、接触点布局太差）。换求解器不会变出物理上不存在的解。QP infeasible 是有用的诊断信号，说明任务超出了当前抓取配置的能力。应对：增大摩擦（换材料/增法向预紧）、改善接触布局（撑开 $\mathbf{r}_i$）、降低运动需求（减小加速度$\to$减小 $\mathbf{w}_{\text{ext}}$）、或加机器人，而非死磕求解器。自检：放松约束（临时调大 $f^{\max}$、$\mu$）看是否变可行，定位是哪类约束被违反。
\end{pitfall}

\begin{practice}
\begin{enumerate}
  \item （$\star\star$）三台机器人正三角形托举一块 $200$ N 的平板（竖直向上），接触均为点接触带摩擦、$\mu=0.5$、法向竖直向上。若用等分策略，每点出力多少？这个等分解满足摩擦锥吗（切向力为 $0$、法向 $66.7$ N）？现在平板要水平加速（需额外水平合力 $60$ N），等分策略如何处理水平力？为什么此时需要 QP？
  \item （$\star\star\star$）用 OSQP 的 Python 绑定实现 \cref{eq:cm-qp-std} 的负载分配 QP，针对练习 1 的三机器人托举场景。加入摩擦金字塔（$k=4$）、单边、力限约束。测试：（a）纯竖直托举时各点力；（b）逐渐增大水平加速度，观察何时 QP 变 infeasible（摩擦锥耗尽）；（c）把 $\mu$ 从 $0.5$ 调到 $0.8$，infeasible 阈值如何变化。
  \item （$\star\star\star$）负载分配 QP 每个控制周期都要解一次（$1$ kHz）。讨论实时性优化：（a）为什么用“热启动（warm start，用上周期解初始化）”能加速 OSQP？（b）若某周期 QP 求解超时，应如何安全降级（提示：回退到上周期解/加权伪逆 $+$ 限幅）？（c）摩擦金字塔边数 $k$ 如何在精度与求解速度间权衡？
\end{enumerate}
\end{practice}

% ════════════════════════════════════════════════════════════════
\section{无源性与稳定性}
\label{sec:cm-passivity}

\paragraph{动机：为什么“每个控制器都稳”不等于“系统稳”。}
\nameref{ch:mr-coopforce} 的“如果跳过本章”场景二：“双机搬运很稳，加第三台就发散。”这个现象击碎了一个常见幻觉——\textbf{子系统稳定 $\not\Rightarrow$ 整体稳定}。三台机器人各自的阻抗控制器单独看都是稳定的（弹簧–阻尼系统），但它们通过物体形成的反馈回路可能注入净能量，使整体发散。为什么？因为稳定性在\emph{互联}下不是简单叠加的——两个稳定系统反馈连接，结果可能不稳（经典控制常识）。多机搬运正是一堆控制器通过物体反馈互联，单独分析每个控制器毫无用处，必须分析互联后的整体。但整体分析太难：维度高（$N$ 台各有动力学）、结构可变（机器人可增减、接触可断合）、含延迟（若有通信）、含时变参数（变阻抗），逐一推 Lyapunov 函数几乎不可能。我们需要一个\textbf{模块化、可组合}的稳定性判据——能对每个子系统单独验证，再断言整体。这把尺就是\textbf{无源性（passivity）}：它的威力在于\emph{无源性在互联下是封闭的}——若干无源子系统按正确拓扑连接，整体仍无源，从而稳定。

\paragraph{反面：只用 Lyapunov 逐系统分析会怎样。}
坚持传统 Lyapunov（为整个 $N$ 机系统构造一个 Lyapunov 函数并证 $\dot V<0$）会遇三个困难。\textbf{其一，维度灾难与重构负担}：$N$ 变了、接触拓扑变了、加了延迟，整个 Lyapunov 函数就得重新构造证明；无源性是模块化的——加一台机器人只需验证它无源，不必重证全局。\textbf{其二，难以处理延迟与时变}：通信延迟使系统成为时滞系统，时变阻抗使系统非自治，二者都让 Lyapunov 函数的构造异常困难；无源性框架有现成工具（散射理论/波变量处理延迟、能量罐处理时变）。\textbf{其三，丢失物理直觉}：Lyapunov 函数常是“凑”出来的数学构造，物理意义不明；无源性的储能函数就是系统的物理能量，$\dot V\le$ 输入功率直接说“系统不造能量”。

\begin{insight}[把“稳定性”重述为“追踪能量的流动”]
无源性把“稳定性”从“求解微分方程的渐近行为”重新表述为“追踪能量的流动”。一个系统稳不稳，等价于问“它会不会凭空生出能量”。这个视角的革命性在于：能量是可叠加、可在子系统间传递、可在互联处守恒的物理量——所以稳定性分析变成了能量记账，而记账是模块化、可组合的。这正是它碾压逐系统 Lyapunov 的根本原因。
\end{insight}

\paragraph{历史：从电路网络到遥操作再到协作。}
无源性源于\textbf{电路网络理论}——无源元件（电阻、电容、电感）不产生能量，无源网络互联仍无源。1960s--70s 这套理论被引入控制（Willems 的耗散系统理论）。1980s--90s，\textbf{遥操作}领域遇到通信延迟导致不稳的难题：Anderson 与 Spong 1989\cite{anderson1989bilateral} 用\textbf{散射理论（scattering theory）}证明延迟通信信道无源的关键条件；Niemeyer 与 Slotine 1991\cite{niemeyer1991stable} 提出\textbf{波变量（wave variables）}，把延迟信道改造成无源，从根本上解决了延迟遥操作的稳定性。2000s，时域无源性方法（Hannaford--Ryu 的 Time-Domain Passivity / 能量罐 energy tank）\cite{hannaford2002time} 提供在线监测并保证无源的工具。这些工具被自然迁移到多机协作搬运——因为“多机器人 $+$ 物体”和“主从遥操作”共享同一个无源性骨架（\cref{sec:cm-leader} 已指出 leader--follower 源于遥操作），综述见 Nu\~no 等\cite{nuno2011passivity}。

\subsection{无源性定义}
\label{sec:cm-passivity-def}
考虑一个系统，输入 $\mathbf{u}$、输出 $\mathbf{y}$（取为“功率共轭对”，如力与速度，$\mathbf{u}^\top\mathbf{y}$ 是功率）。系统\textbf{无源}，若存在下有界的储能函数 $V(\mathbf{x})\ge 0$，使得对所有时间
\begin{equation}
\boxed{\;\dot{V}\le\mathbf{u}^\top\mathbf{y}\quad\Longleftrightarrow\quad V(t)-V(0)\le\int_0^t\mathbf{u}^\top\mathbf{y}\,d\tau.\;}
  \label{eq:cm-passivity}
\end{equation}
含义：系统储能的增量不超过外界注入的能量（功率积分）——系统\emph{不会凭空产生能量}，它至多储存或耗散外界给的能量。若 $\dot V\le\mathbf{u}^\top\mathbf{y}-\delta\|\mathbf{y}\|^2$（$\delta>0$）则\textbf{严格输出无源}（额外耗散）。

\subsection{无源 \texorpdfstring{$\Rightarrow$}{=>} 稳定}
\label{sec:cm-passivity-stable}
取储能函数 $V$ 作 Lyapunov 候选。当外部输入为零（$\mathbf{u}=\mathbf{0}$）时，\cref{eq:cm-passivity} 给出 $\dot V\le 0$——储能不增，系统状态有界、收敛到能量极小点。这就是“无源蕴含（Lyapunov 意义下）稳定”。直观：一个不产能、且能耗散能量的系统，最终会“安定”到低能量状态。

\subsection{互联定理：无源性在反馈/并联下封闭}
\label{sec:cm-passivity-interconnect}
这是无源性最有用的性质。两个无源系统 $\Sigma_1,\Sigma_2$：\emph{负反馈互联}（$\Sigma_1$ 输出经 $\Sigma_2$ 反馈回输入）——整体无源，储能函数取 $V=V_1+V_2$，功率在互联处守恒，整体 $\dot V\le\mathbf{u}_{\text{ext}}^\top\mathbf{y}_{\text{ext}}$；\emph{并联}（同输入、输出相加）——整体无源。推论（本章关键）：\textbf{若每个机器人控制器无源、物体（刚体）无源、环境无源，则整个“机器人–物体–环境”互联系统无源，从而稳定。} 这就是把零散控制律焊成稳定系统的定理。

\paragraph{多机搬运为何天然适合无源性。}
回顾 \cref{sec:cm-decentral} 的信息流 $\mathbf{f}_i\xrightarrow{\mathbf{G}}\mathbf{w}_o\xrightarrow{\text{物体}}\mathbf{v}_o\xrightarrow{\mathbf{G}_j^\top}\dot{\mathbf{x}}_j$。注意 $\mathbf{G}$ 与 $\mathbf{G}^\top$ 的对偶（\cref{sec:cm-grasp} 的功守恒 $\mathbf{f}^\top\dot{\mathbf{X}}=\mathbf{w}_o^\top\mathbf{v}_o$）意味着：\emph{抓取界面本身不产生也不耗散能量}——它是无源的（功率守恒地传递）。物体作为刚体，动能 $\frac12\mathbf{v}_o^\top\mathbf{M}_o\mathbf{v}_o$ 是天然储能函数，物体也无源。所以只要各机器人控制器设计成无源（如标准阻抗律：弹簧势能 $+$ 阻尼耗散），整个系统自动无源稳定。

\begin{insight}[多机搬运的稳定性“免费”继承自三件事]
多机搬运的稳定性“免费”地继承自：（1）抓取界面无源（功守恒，$\mathbf{G}/\mathbf{G}^\top$ 对偶保证）；（2）物体无源（刚体动能作储能）；（3）标准阻抗控制器无源（势能 $+$ 耗散）。三个无源系统互联即稳。麻烦只在\emph{破坏无源性的因素}：变阻抗（注入能量）和通信延迟（产生能量）。本节后半就是修复这两个破坏者。
\end{insight}

\subsection{标准阻抗律为何无源}
\label{sec:cm-passivity-impedance}
验证 \cref{sec:cm-impedance} 物体阻抗（设前馈/重力已补偿，看误差动力学）$\mathbf{M}_d\ddot{\tilde{\mathbf{x}}}+\mathbf{D}_d\dot{\tilde{\mathbf{x}}}+\mathbf{K}_d\tilde{\mathbf{x}}=\mathbf{w}_{\text{env}}$，输入 $\mathbf{w}_{\text{env}}$、输出 $\dot{\tilde{\mathbf{x}}}$。取储能 $V=\frac12\dot{\tilde{\mathbf{x}}}^\top\mathbf{M}_d\dot{\tilde{\mathbf{x}}}+\frac12\tilde{\mathbf{x}}^\top\mathbf{K}_d\tilde{\mathbf{x}}$（动能 $+$ 弹簧势能，$\ge 0$）。求导：
\begin{equation}
  \dot V=\dot{\tilde{\mathbf{x}}}^\top\mathbf{M}_d\ddot{\tilde{\mathbf{x}}}+\dot{\tilde{\mathbf{x}}}^\top\mathbf{K}_d\tilde{\mathbf{x}}
  =\dot{\tilde{\mathbf{x}}}^\top(\mathbf{w}_{\text{env}}-\mathbf{D}_d\dot{\tilde{\mathbf{x}}}-\mathbf{K}_d\tilde{\mathbf{x}})+\dot{\tilde{\mathbf{x}}}^\top\mathbf{K}_d\tilde{\mathbf{x}}
  =\dot{\tilde{\mathbf{x}}}^\top\mathbf{w}_{\text{env}}-\dot{\tilde{\mathbf{x}}}^\top\mathbf{D}_d\dot{\tilde{\mathbf{x}}}.
  \label{eq:cm-imp-passive}
\end{equation}
由于 $\mathbf{D}_d\succ 0$，$\dot{\tilde{\mathbf{x}}}^\top\mathbf{D}_d\dot{\tilde{\mathbf{x}}}\ge 0$，故 $\dot V\le\dot{\tilde{\mathbf{x}}}^\top\mathbf{w}_{\text{env}}=\mathbf{u}^\top\mathbf{y}$。\textbf{满足无源性 \cref{eq:cm-passivity}，且严格输出无源（阻尼项额外耗散）。} 这就是固定参数阻抗律稳定的能量证明——只要 $\mathbf{M}_d,\mathbf{D}_d,\mathbf{K}_d$ 是常正定矩阵。

\subsection{破坏者一：变阻抗与能量罐}
\label{sec:cm-passivity-tank}
\paragraph{为什么变阻抗危险。}
\cref{sec:cm-impedance} 的思维陷阱与单臂力控都提到：在线改变刚度 $\mathbf{K}_d$ 会注入能量。看上面的证明——若 $\mathbf{K}_d=\mathbf{K}_d(t)$ 时变，储能 $V=\dots+\frac12\tilde{\mathbf{x}}^\top\mathbf{K}_d(t)\tilde{\mathbf{x}}$ 求导多出一项 $\frac12\tilde{\mathbf{x}}^\top\dot{\mathbf{K}}_d\tilde{\mathbf{x}}$：
\begin{equation}
  \dot V=\dot{\tilde{\mathbf{x}}}^\top\mathbf{w}_{\text{env}}-\dot{\tilde{\mathbf{x}}}^\top\mathbf{D}_d\dot{\tilde{\mathbf{x}}}+\tfrac12\tilde{\mathbf{x}}^\top\dot{\mathbf{K}}_d\,\tilde{\mathbf{x}}.
  \label{eq:cm-varimp}
\end{equation}
当 $\dot{\mathbf{K}}_d\succ 0$（刚度增大）时，$\frac12\tilde{\mathbf{x}}^\top\dot{\mathbf{K}}_d\tilde{\mathbf{x}}>0$ \textbf{凭空注入能量}，可能使 $\dot V>\mathbf{u}^\top\mathbf{y}$——无源性被破坏，系统可能失稳。物理上：增大弹簧刚度（在弹簧已被拉伸时）相当于“给弹簧充能”，这能量来自控制器，是主动注入。多机搬运里每台机器人若用变阻抗（如接近时硬、接触时软）都可能成为产能源。

\paragraph{能量罐：给能量注入装一个“预算账户”。}
能量罐（energy tank）的思想：设一个虚拟能量储罐 $T_{\text{tank}}\ge 0$，\emph{所有“产能”的动作（如增大刚度）必须从罐里取能}，罐空了就不许产能（冻结刚度变化或限制其方向）；罐的能量来自系统的耗散（阻尼耗散的能量被回收进罐）。这样系统对外永远无源——产能动作用的是之前耗散攒下的能量，不是凭空创造。能量罐动力学（示意）：
\begin{equation}
  \dot T_{\text{tank}}=\underbrace{\dot{\tilde{\mathbf{x}}}^\top\mathbf{D}_d\dot{\tilde{\mathbf{x}}}}_{\text{回收耗散}}-\underbrace{P_{\text{inject}}}_{\text{产能动作取能}},\qquad T_{\text{tank}}\ge T_{\min}>0.
  \label{eq:cm-tank}
\end{equation}
放行规则：当要做产能动作（$P_{\text{inject}}>0$，如增刚度）时，仅当 $T_{\text{tank}}>T_{\min}$ 才放行，否则按比例缩减或禁止。整体储能取 $V_{\text{total}}=V+T_{\text{tank}}$，可证 $\dot V_{\text{total}}\le\mathbf{u}^\top\mathbf{y}$——\textbf{带能量罐的变阻抗系统重新无源}。这是单臂能量罐在多机的直接推广：每台机器人挂一个能量罐，或物体级挂一个共享罐。

\subsection{破坏者二：通信延迟与波变量}
\label{sec:cm-passivity-wave}
\paragraph{为什么延迟产生能量。}
若多机搬运确实需要通信（如交换状态做协调），通信延迟 $T$ 会让“延迟信道”成为产能元件。直观：你在 $t$ 时刻发出的力/速度信息，对方在 $t+T$ 收到，这个时间错位使“发送端做的功”和“接收端收到的功”不匹配，差额可正可负——当为正时信道凭空“创造”了能量，破坏无源性，导致延迟越大越易失稳。这正是早期遥操作“延迟一大就发散”的根源。

\paragraph{波变量：把延迟信道改造成无源。}
波变量（Niemeyer--Slotine\cite{niemeyer1991stable}）的精妙：不直接传“力”和“速度”，而传它们的\emph{线性组合}——波变量。定义（一维示意，$b$ 为波阻抗）：
\begin{equation}
  u=\frac{1}{\sqrt{2b}}(F+b\,\dot x),\qquad v=\frac{1}{\sqrt{2b}}(F-b\,\dot x),
  \label{eq:cm-wave}
\end{equation}
$u$ 是前向波、$v$ 是返回波。关键性质：用波变量传输时，\textbf{无论延迟 $T$ 多大，信道储能 $\ge 0$（无源）}。证明核心：信道功率 $P=\frac12(\|u\|^2-\|v\|^2)$，延迟只是把 $u(t)$ 变成 $u(t-T)$、$v(t)$ 变成 $v(t-T)$，可证信道在 $[0,t]$ 净吸收能量 $\ge 0$——延迟信道用波变量传输后变成无源的“传输线”，不再产能。代价：波变量引入“波反射”等动态，需调 $b$ 折中响应与稳定。

\begin{note}[类比的边界：波变量 vs 长管中的声波]
波变量之于延迟通信，像声音在长管中以波的形式传播。\emph{像的部分}：信息以“波”形式传递、有传播延迟、用特征阻抗 $b$ 描述传输线、不会凭空增能（无源传输线）。\emph{不像的部分}：物理声波是被动现象，波变量是主动编码控制信号的数学变换；且真实通信是离散采样的，波变量需配合采样无源化处理（否则离散化又会引入能量）。把“声波天然无源”照搬，会忽略离散实现时的额外无源化需求。
\end{note}

\paragraph{多机搬运里何时需要波变量。}
纯去中心化（\cref{sec:cm-decentral}，零通信）不需要波变量——它根本不通过网络传信息，靠物体（无源的物理信道）。波变量只在\emph{确实需要显式通信协调}且延迟不可忽略时才需要（如机器人相距远、用无线网交换状态）。这给了一个清晰的工程判断：能用物体隐式通信（\cref{sec:cm-decentral}）就用，物理信道天然无源、零延迟；必须显式通信且有延迟时，用波变量把网络信道无源化。

\paragraph{代码实现要点（能量罐放行阀）。}
能量罐每周期：先回收阻尼耗散（$T\mathrel{+}=$ dissipated，限幅到 $T_{\max}$ 防止无限积累），耗能动作（$P_{\text{inject}}\le 0$）直接放行；产能动作（$P_{\text{inject}}>0$）只放行 $\min(P_{\text{inject}},\max(0,T-T_{\min}))$ 并从罐扣除，调用方据此缩放刚度变化率。两处常见错误：（1）不设 $T_{\min}$，允许罐耗尽到 $0$ 仍放行产能动作 $\to$ 实际注入了不存在的能量 $\to$ 无源性破坏，必须保留 $T_{\min}>0$ 缓冲；（2）不限幅 $T_{\max}$，长期耗散使罐无限增大 $\to$ 攒了巨量“能量预算”，某刻可一次性大注入、形同没有保护。

\begin{insight}[无源性如何指导你调多机搬运的增益]
无源性不只是事后判稳，更是设计指南：（1）阻尼 $\mathbf{D}_d$ 必须正定且足够大——它是系统唯一的耗散源，太小则接近无源边界、易振荡；（2）刚度变化要挂能量罐，否则变阻抗会偷偷产能；（3）有通信延迟就用波变量，别指望调增益硬扛延迟。“加第三台就发散”，用无源性诊断就是：检查新加的机器人控制器是否无源（阻尼够不够、有没有不受控的产能项）、互联是否引入了产能通道。这把抽象的能量记账变成了具体的调参清单。
\end{insight}

\begin{pitfall}[思维陷阱：以为各子系统稳定则整体稳定]
“每台机器人的阻抗控制器都是稳定的弹簧–阻尼，合起来当然稳”——稳定性在反馈互联下不叠加，两个稳定系统反馈连接可能不稳。多机通过物体强反馈互联，必须分析互联后的整体。正确的模块化判据是无源性（不是稳定性）：各子系统无源 $+$ 正确互联 $\Rightarrow$ 整体无源 $\Rightarrow$ 稳定。无源性可叠加，稳定性不可。验证“无源”而非“稳定”作为模块化判据，问每个子系统“会不会产能”而非“自己稳不稳”。自检：为每个控制器写储能函数，验 $\dot V\le\mathbf{u}^\top\mathbf{y}$；若某控制器产能（如未保护的变阻抗），它就是整体失稳的嫌疑源。
\end{pitfall}

\begin{pitfall}[概念误区：把无源性等同于稳定性]
“无源就是稳定，两个词换着用”——无源 $\Rightarrow$（Lyapunov 意义）稳定，但反之不真：存在稳定却非无源的系统。无源是比稳定更强、且可组合的性质，我们用无源正是因为它能模块化组合（互联封闭），稳定性本身不能；且无源是关于“能量”的输入输出性质，稳定是关于“状态收敛”的内部性质，概念层次不同。混淆二者会让你误以为“证了稳定就等于证了可安全互联”，但只有无源才保证互联后仍稳。延伸：严格输出无源 $+$ 可检测 $\Rightarrow$ 渐近稳定，这类定理把无源和稳定精确联系起来，但二者不是等号。
\end{pitfall}

\begin{pitfall}[编程陷阱：能量罐回收的“耗散能量”算错符号或来源]
把“环境注入的能量”也当耗散回收进罐，或把阻尼耗散符号算反（本应 $\ge 0$ 却出现负值）。现象：罐能量虚高，系统获得了不属于自己耗散的“能量预算”，产能动作被过度放行，无源性保护形同虚设，变阻抗仍致失稳。根因：能量罐只能回收系统自身耗散（阻尼项 $\dot{\tilde{\mathbf{x}}}^\top\mathbf{D}_d\dot{\tilde{\mathbf{x}}}\ge 0$）的能量，不能把外界注入或测量噪声当能量收。对策：严格只回收阻尼耗散（恒 $\ge 0$）、验证回收量非负、环境交互能量不进罐。自检：系统静止（$\dot{\tilde{\mathbf{x}}}=\mathbf{0}$）时回收应为 $0$，罐能量不应在静止时增长。
\end{pitfall}

\begin{practice}
\begin{enumerate}
  \item （$\star\star\star$）完整证明：固定参数的物体阻抗误差动力学 $\mathbf{M}_d\ddot{\tilde{\mathbf{x}}}+\mathbf{D}_d\dot{\tilde{\mathbf{x}}}+\mathbf{K}_d\tilde{\mathbf{x}}=\mathbf{w}_{\text{env}}$（$\mathbf{M}_d,\mathbf{D}_d,\mathbf{K}_d$ 常正定）以 $\mathbf{w}_{\text{env}}$ 为输入、$\dot{\tilde{\mathbf{x}}}$ 为输出是无源的。然后证明：若 $\mathbf{K}_d=\mathbf{K}_d(t)$ 单调增大，无源性可能被破坏，并指出被注入能量的具体项。最后说明能量罐如何修复（写出带罐的总储能函数）。
  \item （$\star\star\star$）“加第三台机器人就发散。”用无源性框架给出至少三个可能的“产能源”诊断假设（如：第三台控制器阻尼不足接近无源边界、用了未保护的变阻抗、它与其他机器人间存在通信延迟未做波变量处理）。对每个假设给出一个可操作的检验方法和对应修复措施。
  \item （$\star\star\star\star$，跨章综合）论述无源性这把“尺”如何贯穿三个场景：单臂变阻抗（能量罐）、主从遥操作（波变量处理延迟）、多机协同搬运（互联封闭定理）。三者的“储能函数”“破坏无源的因素”“修复手段”分别是什么？为什么同一套无源性理论能统一这三个看似不同的问题——它们共享的物理本质（能量流动）是什么？
\end{enumerate}
\end{practice}

% ════════════════════════════════════════════════════════════════
\section{前沿工作与综合}
\label{sec:cm-frontier}

前七节建立了完整的理论与方法工具箱。这一节用前面的概念框架解读三个代表性前沿工作，看经典理论如何在最新系统里活着，并给出一个综合实战的累积项目骨架。

\subsection{前沿一：多足绳索搬运（Sombolestan \& Nguyen 2022）}
\label{sec:cm-frontier-cable}
\textbf{问题与方案}：多台四足机器人通过\emph{绳索}牵引一个负载协同搬运\cite{sombolestan2022collaborative}。\rebuilt{引用张冠李戴：标题《...a Cable-towed Load by Multiple Quadrupedal Robots》实为 Yang 等(UC Berkeley)、发表于 RA-L 7(4):10041--10048(2022)（非 IROS、非 Sombolestan \& Nguyen），且其核心是\emph{级联轨迹优化}（利用绳松/紧切换过窄道）而非分布式力控。Sombolestan \& Nguyen(USC)确有相关四足协同操作工作（IROS 2023 刚体自适应力控），但与本标题非同一篇。键 sombolestan2022collaborative 的作者/venue/内容均待修正；下文"分布式力控"解读更贴近 Sombolestan \& Nguyen 的真实工作，而非所引标题之文。}绳索是关键约束——它\textbf{只能拉、不能推}（单边约束的极端形式），且几乎不能传力矩。\textbf{用本章框架解读}（以"多足绳索/单边约束搬运"为一般情形，不绑定上述误引的具体论文）：
\begin{itemize}
  \item \textbf{抓取矩阵}：每根绳索的接触力 $\mathbf{f}_i$ 沿绳方向、只能为拉力（$\mathbf{f}_i$ 在绳方向投影 $\ge 0$）；$\mathbf{G}$ 把各绳张力映射到负载旋量，结构同 \cref{sec:cm-grasp}，但接触力被严格限制在“沿绳、且为拉”的一维射线上。
  \item \textbf{负载分配}：变成 \cref{sec:cm-qp} 的 QP，约束特别强——每根绳张力 $\ge 0$（单边）、$\le$ 绳强度上限，摩擦锥退化为“张力非负”。这是“单边约束”最纯粹的体现。
  \item \textbf{架构}：分布式力控——每台四足本地控制自己的绳张力，靠负载运动隐式协调（\cref{sec:cm-decentral} 思想）。
  \item \textbf{稳定性}：绳的柔性引入额外动态，无源性分析（\cref{sec:cm-passivity}）需考虑绳作为弹性储能元件。
\end{itemize}

\begin{insight}[绳索搬运是多机搬运的“约束极限案例”]
绳索搬运把“接触力可任意方向”收紧到“只能沿一根线拉”。它逼着你直面 \cref{sec:cm-qp} 的单边约束和 \cref{sec:cm-internal} 的内力含义（多根绳之间的张力分配就是内力调节：绳越紧内力越大、越稳但越费力）。理解了绳索搬运，你对一般接触的内力/可行性约束会有更深的体会。
\end{insight}

\subsection{前沿二：被动臂搬运（PACC, Turrisi 2024）}
\label{sec:cm-frontier-pacc}
\textbf{问题与方案}：高负载协同搬运中，机器人间的耦合（一台动，通过刚性物体扰动其他台）让控制极难。PACC 的巧思\cite{turrisi2024pacc}：不在四足足端硬抓物体（高耦合），而是在每台四足背上装一根\textbf{被动臂（passive arm，无电机）}，臂端与负载柔顺连接\pz{原文该被动臂为 3 个转动关节(yaw-pitch-pitch)$+$弹簧/阻尼实现内在阻抗，源稿"万向节"系臆造}。\textbf{用本章框架解读}：
\begin{itemize}
  \item \textbf{耦合解耦}：被动臂的被动关节像“机械低通滤波器”，\emph{吸收}机器人与负载间的高频交互力。这相当于在 \cref{sec:cm-grasp} 的刚性闭链里插入了柔性环节——闭链不再刚性，内力被被动臂的形变吸收，不会瞬间传遍全系统。
  \item \textbf{简化控制}：因为被动臂吸收了耦合，每台四足可跑自己的 MPC（如 OCS2），几乎不必显式处理多机耦合——这是用\emph{硬件（被动臂）换软件复杂度}的工程权衡。\pz{"标准 MPC"措辞不准：原文 MPC 是\emph{增广}的——纳入了臂动力学近似与外部力估计，并改了动态稳定判据，此恰为该文核心贡献，非现成标准 MPC。}
  \item \textbf{代价}：需额外硬件（被动臂 $+$ 万向节），最大负载受臂结构强度限制。
\end{itemize}

\begin{note}[类比的边界：被动臂 vs 汽车悬架]
被动臂之于刚性抓持，像汽车悬架之于硬连接车轴。\emph{像的部分}：都用被动柔性元件吸收冲击/耦合，让上游（车身/四足控制器）不必处理高频扰动。\emph{不像的部分}：悬架主要隔振，被动臂还改变了闭链的力传递结构（把内力变成被动臂的形变）；且被动臂的引入改变了 \cref{sec:cm-grasp} 抓取矩阵的刚性假设。把“悬架 $=$ 纯隔振”照搬会低估被动臂对内力空间的重构作用。
\end{note}

\subsection{前沿三：零通信学习搬运（decPLM, Pandit 2025）}
\label{sec:cm-frontier-decplm}
\textbf{问题与方案}：$2$--$10$ 台 Go2 四足\textbf{零通信}协同搬运，用强化学习（IsaacLab 训练）学出 \textbf{pinch-lift-move}（捏紧$\to$抬起$\to$移动）策略\cite{pandit2025decplm}。每台 Go2 只用本体感知（自身姿态、足端接触/力），不与其他 Go2 通信。\textbf{用本章框架解读}：
\begin{itemize}
  \item \textbf{去中心化 $+$ 隐式通信}：这是 \cref{sec:cm-decentral} 的纯粹实现——同构策略、零通信、靠物体运动隐式协调，协调从物理耦合中涌现。
  \item \textbf{pinch $=$ 内力调节}（解读性映射）：捏紧阶段\emph{可解读为}在 $\ker\mathbf{G}$ 里增大内力（夹紧防滑，\cref{sec:cm-internal}）；lift/move \emph{可解读为}施加外力旋量（\cref{sec:cm-impedance} 物体运动）。\pz{此 $\ker\mathbf{G}$/内外力映射是本书用经典框架对该 RL 工作的\emph{解读}，非 decPLM 原文主张——原文三阶段为训练课程、动作空间是速度/位置命令而非力旋量，全文不含抓取矩阵/内力分解；勿当作其自身的力空间分解。}
  \item \textbf{学习替代显式建模}：经典方法（\cref{sec:cm-impedance}）需要物体惯量 $\mathbf{M}_o$，decPLM 用 RL 免模型——策略从经验中隐式习得如何应对未知负载，呼应脉络主线“从精确模型走向自适应/学习”的终点。
\end{itemize}

\begin{insight}[decPLM 用神经网络隐式学会了本章显式推导的东西]
decPLM 不是抛弃了本章理论，而是用神经网络\emph{隐式地学会了}本章显式推导的东西——它的策略内部一定编码了某种“内力/外力分解”和“隐式协调”，只是没写成 $\mathbf{G}^{+}\mathbf{w}+\mathbf{E}\boldsymbol{\lambda}$ 的解析形式。理解本章理论，你才能看懂 decPLM 学到了什么、它的 pinch-lift-move 为何对应内力调节与外力施加，以及它的局限（学习策略缺乏无源性保证，\cref{sec:cm-passivity} 的稳定性证明对它不直接适用，这正是学习型搬运的开放问题）。
\end{insight}

\begin{table}[H]\centering\small
\caption{前沿全景与开放问题（按源稿重述，待核）}
\label{tab:cm-frontier}
\begin{tabular}{p{2.6cm}p{2.0cm}p{2.2cm}p{1.8cm}p{2.0cm}p{2.4cm}}
\toprule
工作 & 架构 & 接触 & 建模 & 本章对应 & 开放问题 \\
\midrule
Sombolestan 2022 & 分布式 & 绳索（单边拉） & 模型 & \cref{sec:cm-decentral}$+$\cref{sec:cm-qp} 单边 QP & 绳柔性下的稳定性\pz{引用张冠李戴，见 \cref{sec:cm-frontier-cable} 注} \\
PACC 2024 & 各自增广 MPC & 被动臂(3R$+$弹簧阻尼) & 模型（解耦） & \cref{sec:cm-grasp} 柔性闭链 & 被动臂参数设计、负载上限\pz{源稿"万向节/标准 MPC"有误，见 \cref{sec:cm-frontier-pacc} 注} \\
decPLM 2025 & 去中心化学习 & pinch（足） & 免模型 RL & \cref{sec:cm-decentral}$+$\cref{sec:cm-internal} 内力 & 学习策略的稳定性保证、泛化 \\
Verginis 2019 & 分布式自适应 & 刚性 & 自适应（免惯量） & \cref{sec:cm-impedance}$+$\cref{sec:cm-passivity} & 动态（非准静态）搬运\pz{发表 venue 为 IEEE T-CNS(2023) 非 RA-L} \\
\bottomrule
\end{tabular}
\end{table}

\noindent\textbf{核心开放问题}：（1）\textbf{动态搬运的内力优化}——负载惯性效应显著（如四足跑步时搬运）时如何实时优化内力？现有多数方法（Verginis）限准静态。（2）\textbf{学习型搬运的稳定性保证}——decPLM 等 RL 策略缺乏 \cref{sec:cm-passivity} 的无源性证明，如何把无源性约束嵌入学习？（3）\textbf{柔性物体搬运}——本章假设刚体，真实物体（布料、长杆、液体容器）的柔性使闭链约束变成偏微分方程，理论尚不成熟。

\subsection{累积项目：Mini-Cooperative-Transport 本章模块}
\label{sec:cm-frontier-project}
本章在累积项目（从零搭建多机协作系统）中新增\textbf{协同搬运与力控模块}，复用前几章的通信图建模、共识协调、任务分配/MAPF、分布式 MPC 编队。实战目标：在仿真（如 MuJoCo）中搭建双 Go2 $+$ 刚性杆场景，实现协同搬运。分阶段递进（渐隐示例法：从给定到独立）：
\begin{itemize}
  \item \textbf{Phase 1（完全跟随）}：给定场景（两 Go2 $+$ 刚性杆，杆两端与各 Go2 固连），每帧取两 Go2 抓持点世界坐标与杆质心，调用抓取矩阵组装算当前 $\mathbf{G}$，验证 $\rank(\mathbf{G})$ 与内力维度（双机点接触 $\Rightarrow$ 内力维度 $0$）。
  \item \textbf{Phase 2（填空）}：给定物体期望轨迹（杆从地面抬到目标高度），用物体阻抗 \crefrange{eq:cm-obj-imp}{eq:cm-imp-compose} 算所需旋量、用负载分配 QP \cref{eq:cm-qp-std} 分给两 Go2。关键步骤留空：杆位姿误差、速度误差、内力意图 $\boldsymbol{\lambda}_{\text{int}}$ 的取法、\textbf{重力旋量务必补偿}。
  \item \textbf{Phase 3（主导）}：改写为去中心化（\cref{sec:cm-decentral}）——每个 Go2 跑同一本地控制律、仅用本地观测。测试鲁棒性：（a）一个 Go2 注入定位噪声，观察隐式协调是否仍收敛；（b）给杆施加外部冲击，验证物体阻抗的柔顺响应；（c）用能量罐保护一个变阻抗策略，对比有/无能量罐的稳定性。
  \item \textbf{Phase 4（独立）}：扩展为三机（内力维度从 $0$ 变 $3$，需主动调内力），或改为绳索牵引（单边约束，QP 必须处理）。
\end{itemize}
\textbf{验收标准}：杆被平稳抬起并跟踪目标轨迹（位置误差 $<5$ cm）、两 Go2 接触力之和竖直分量 $\approx$ 杆重、内力维持合理范围（不挤爆、不脱手）、去中心化版本在定位噪声下仍收敛。

\begin{practice}
\begin{enumerate}
  \item （$\star\star\star$）精读 decPLM 的方法部分，用本章语言标注：（a）它的“去中心化”对应 \cref{sec:cm-decentral} 的哪些要素？（b）pinch-lift-move 三阶段分别对应内力/外力的什么操作？（c）它如何绕开物体惯量未知（对比 \cref{sec:cm-impedance} 基于模型 vs Verginis 自适应 vs decPLM 学习）？（d）它缺少 \cref{sec:cm-passivity} 的什么保证？
  \item （$\star\star\star$）完成累积项目 Phase 1--2：在仿真搭双 Go2 $+$ 刚性杆场景，实现物体阻抗 $+$ 伪逆力分配，把杆从地面抬到 $0.5$ m 高。报告：物体跟踪误差曲线、两 Go2 接触力曲线、验证接触力之和竖直分量 $\approx$ 杆重。
  \item （$\star\star\star\star$，跨章综合）设计一个完整的多机搬运任务管线，串起本 Part 全部知识：任务分配决定“哪几台机器人去搬哪个物体”、MAPF/编队 MPC 让它们无碰撞地移动到物体旁、本章 \crefrange{sec:cm-grasp}{sec:cm-passivity} 完成抓取与协同搬运。画出完整数据流图，标注每步用到的章节、关键算法与接口。讨论：从“移动到物体旁”切换到“协同搬运”的瞬间，控制架构如何平滑过渡（编队 MPC $\to$ 闭链力控）？这个切换点最容易出什么问题？
\end{enumerate}
\end{practice}

% ════════════════════════════════════════════════════════════════
\section*{本章常见误解汇总}
\addcontentsline{toc}{section}{本章常见误解汇总}

\begin{table}[H]\centering\small
\caption{协同搬运与多机力控常见误解}
\label{tab:cm-misconceptions}
\begin{tabular}{p{4.6cm}p{6.8cm}c}
\toprule
误解 & 正确理解 & 出处 \\
\midrule
抓取矩阵是多机搬运的专用工具矩阵 & 它是闭链约束的数学化身，与单臂雅可比 $\mathbf{J}$ 同构（力映射 $\mathbf{G}$、速度对偶 $\mathbf{G}^\top$），零空间住着内力 & \cref{sec:cm-grasp} \\
两台机器人一起搬就一定有内力要控制 & 内力维度 $=cN-6$：点接触双机（$c{=}3$）内力维度为 $0$；刚性抓持（$c{=}6$）才有 $6$ 维内力 & \cref{sec:cm-grasp} \\
Moore--Penrose 伪逆解自动给零内力/非挤压分配 & 最小范数 $\ne$ 非挤压；真实搬运需非零夹紧力，内力须显式设计 & \cref{sec:cm-internal} \\
内力是应当消除的副作用，越小越好 & 内力是抓持必需品（防滑），目标是“调到恰当值”（防滑下界、不损坏上界） & \cref{sec:cm-internal} \\
follower 必须知道物体绝对轨迹才能配合 & follower 只需物体当前运动方向——意图经物体运动隐式传递，不需轨迹广播 & \cref{sec:cm-leader} \\
leader 出多大力物体就被多大力驱动 & 力放大架构里 follower 力之和可远大于 leader 力；leader 定方向、follower 提供动力 & \cref{sec:cm-leader} \\
去中心化 $=$ 机器人各干各的、互不影响 & 去中心化机器人通过物体强耦合，协调从物理耦合涌现；“独立”的只是计算 & \cref{sec:cm-decentral} \\
去中心化总优于集中式 & 去中心化用“放弃全局信息”换可扩展与鲁棒，代价是精度；架构是权衡不是优劣排序 & \cref{sec:cm-decentral} \\
多机阻抗 $=$ $N$ 个末端各跑一套阻抗 & 物体阻抗在 $6$ 维物体空间统一定义（运动/内力解耦），$N$ 个末端独立阻抗会过定义 & \cref{sec:cm-impedance} \\
加大物体刚度 $\mathbf{K}_d$ 总能提升跟踪精度 & $\mathbf{K}_d$ 过大牺牲柔顺、放大内力、逼近稳定边界；先查前馈/重力补偿/模型 & \cref{sec:cm-impedance} \\
加权伪逆能满足摩擦/单边约束 & 加权伪逆只改“选哪个最小范数解”，仍无约束，照样可能要求拉接触或超摩擦锥 & \cref{sec:cm-qp} \\
QP 无解是求解器问题 & QP infeasible 常是物理事实（物体太重/摩擦太小/布局太差），是有用的诊断信号 & \cref{sec:cm-qp} \\
各子系统稳定则整体稳定 & 稳定性在反馈互联下不叠加；模块化判据是无源性（可叠加），不是稳定性 & \cref{sec:cm-passivity} \\
无源性 $=$ 稳定性 & 无源 $\Rightarrow$ 稳定，反之不真；无源是更强且可组合的性质，这正是用它的原因 & \cref{sec:cm-passivity} \\
\bottomrule
\end{tabular}
\end{table}

% ════════════════════════════════════════════════════════════════
\section*{本章小结}
\addcontentsline{toc}{section}{本章小结}

本章从一个朴素观察出发——多台机器人抓住同一物体就被焊成闭合运动链——系统地建立了协同搬运与多机力控的完整知识体系。\textbf{主线回顾}：（1）\textbf{闭链约束 $\to$ 抓取矩阵}（\cref{sec:cm-grasp}）：闭链耦合的数学化身是抓取矩阵 $\mathbf{G}$，力映射 $\mathbf{w}_o=\mathbf{G}\mathbf{f}$ 与速度对偶 $\dot{\mathbf{X}}=\mathbf{G}^\top\mathbf{v}_o$ 由功守恒联系，零空间维度 $cN-6$ 量化内力自由度。（2）\textbf{内力/外力分解}（\cref{sec:cm-internal}）：接触力正交分解为运动力（$\rowsp\mathbf{G}$，搬物体）与内力（$\ker\mathbf{G}$，夹物体），$\mathbf{f}=\mathbf{G}^{+}\mathbf{w}_{\text{ext}}+\mathbf{E}\boldsymbol{\lambda}$，二者独立设计。（3）\textbf{谁来决策}（\cref{sec:cm-leader,sec:cm-decentral}）：leader--follower（主从、意图估计、力放大）与去中心化（同构、隐式通信、物体即共享黑板）两条路线，构成“全集中—主从—去中心化”权衡谱。（4）\textbf{用什么行为搬}（\cref{sec:cm-impedance}）：物体阻抗把期望柔顺定义在物体质心（而非各末端），运动阻抗与内力阻抗在正交子空间解耦。（5）\textbf{力怎么分}（\cref{sec:cm-qp}）：负载分配从伪逆升级为带摩擦锥/单边/力限的 QP，保证物理可行。（6）\textbf{稳不稳}（\cref{sec:cm-passivity}）：无源性作为模块化判据（互联封闭），能量罐修复变阻抗、波变量修复延迟。（7）\textbf{前沿与综合}（\cref{sec:cm-frontier}）：绳索/被动臂/零通信学习三个前沿用经典框架解读。

\begin{insight}[贯穿全章的本质洞察]
多机搬运的一切——耦合、内力、协调、稳定——都源自“被搬物体把机器人焊成闭链”这一事实。物体既是约束（闭链）、又是隐式通信信道（零通信协作的根基）、还是无源的能量传递媒介（稳定性的来源）。抓取矩阵 $\mathbf{G}$ 及其零空间是描述这一切的统一语言，而它与单臂雅可比、多指抓取共享同一数学结构。
\end{insight}

\subsection*{术语速查表}
\begin{table}[H]\centering\small
\caption{协同搬运与多机力控术语速查（中/英）}
\label{tab:cm-glossary}
\begin{tabular}{p{5.2cm}p{8.0cm}}
\toprule
术语（中 / 英） & 一句话定义 \\
\midrule
抓取矩阵（Grasp Matrix, $\mathbf{G}$） & 把各接触力映射到物体合力旋量的 $6\times cN$ 矩阵，$\mathbf{w}_o=\mathbf{G}\mathbf{f}$ \\
闭合运动链（Closed Kinematic Chain） & 多机器人 $+$ 刚体 $+$ 共同基座构成的力/运动相互约束的环 \\
力旋量（Wrench, $\mathbf{w}$） & 力 $+$ 力矩的 $6$ 维组合，$[\mathbf{F}^\top,\boldsymbol{\tau}^\top]^\top$ \\
内力（Internal Force, $\mathbf{f}_{\text{int}}$） & 位于 $\ker\mathbf{G}$ 的接触力，不产生物体运动，维持抓持（夹紧/张紧） \\
运动力/外力（Motion/External Force） & 位于 $\rowsp\mathbf{G}$ 的接触力，产生期望物体旋量，$\mathbf{G}^{+}\mathbf{w}_{\text{ext}}$ \\
零空间投影（Null-space Projection, $\mathbf{I}-\mathbf{G}^{+}\mathbf{G}$） & 把任意接触力投到 $\ker\mathbf{G}$（内力子空间）的正交投影 \\
非挤压伪逆（Nonsqueezing Pseudoinverse） & Walker 构造的特殊加权伪逆，使力分配不在物体内产生挤压 \\
leader--follower（主从架构） & 一台 leader 知目标用阻抗领航，follower 靠物体运动估计意图配合 \\
去中心化（Decentralized） & 所有 agent 同构、仅用本地测量、不通信，靠物体隐式协调 \\
隐式通信（Implicit Communication） & 通过被搬物体的运动传递意图/状态，无显式网络通信 \\
物体阻抗（Object Impedance） & 把期望阻抗（$\mathbf{M}_d,\mathbf{D}_d,\mathbf{K}_d$）加在物体质心，使物体整体柔顺 \\
负载分配（Load Distribution） & 把物体所需外力旋量分配给各机器人接触力的问题 \\
摩擦锥/摩擦金字塔（Friction Cone/Pyramid） & 切向力 $\le\mu\times$ 法向力 的约束；线性化为多面锥便于 QP \\
单边接触（Unilateral Contact） & 接触只能推或只能拉（如脚只能推、绳只能拉），法向力非负 \\
无源性（Passivity） & 存在储能 $V\ge 0$ 使 $\dot V\le\mathbf{u}^\top\mathbf{y}$，系统不凭空产能 \\
能量罐（Energy Tank） & 给产能动作（如变阻抗）设的能量预算账户，由耗散回收充能 \\
波变量（Wave Variables） & 力与速度的线性组合，使延迟通信信道无源化 \\
波阻抗（Wave Impedance, $b$） & 波变量定义中的特征阻抗参数，折中响应与稳定 \\
\bottomrule
\end{tabular}
\end{table}

\subsection*{核心公式速查}
\begin{table}[H]\centering\small
\caption{协同搬运与多机力控核心公式速查}
\label{tab:cm-formulas}
\begin{tabular}{p{4.0cm}p{7.4cm}c}
\toprule
公式 & 表达式 & 出处 \\
\midrule
抓取矩阵力映射 & $\mathbf{w}_o=\mathbf{G}\mathbf{f}$，$\mathbf{G}_i=[\mathbf{I}_3;\skewm{\mathbf{r}}_i]$ & \cref{eq:cm-grasp} \\
速度对偶 & $\dot{\mathbf{X}}=\mathbf{G}^\top\mathbf{v}_o$ & \cref{eq:cm-vel-dual} \\
内力维度 & $\dim\ker\mathbf{G}=cN-6$ & \cref{eq:cm-null-dim} \\
内力/外力分解 & $\mathbf{f}=\mathbf{G}^{+}\mathbf{w}_{\text{ext}}+(\mathbf{I}-\mathbf{G}^{+}\mathbf{G})\boldsymbol{\eta}=\mathbf{G}^{+}\mathbf{w}_{\text{ext}}+\mathbf{E}\boldsymbol{\lambda}$ & \cref{eq:cm-decomp} \\
右伪逆 & $\mathbf{G}^{+}=\mathbf{G}^\top(\mathbf{G}\mathbf{G}^\top)^{-1}$ & \cref{sec:cm-internal} \\
物体阻抗 & $\mathbf{M}_d\ddot{\tilde{\mathbf{x}}}+\mathbf{D}_d\dot{\tilde{\mathbf{x}}}+\mathbf{K}_d\tilde{\mathbf{x}}=\mathbf{w}_{\text{env}}$ & \cref{eq:cm-obj-imp} \\
所需合旋量 & $\mathbf{w}_{\text{ext}}=\mathbf{M}_o\ddot{\mathbf{x}}^{\text{des}}+\mathbf{C}_o-\mathbf{w}_{\text{grav}}-\mathbf{w}_{\text{env}}$ & \cref{eq:cm-wrench-needed} \\
负载分配 QP & $\min\tfrac12\mathbf{f}^\top\mathbf{W}\mathbf{f}+\mathbf{g}^\top\mathbf{f}$ s.t. $\mathbf{G}\mathbf{f}=\mathbf{w}_{\text{ext}},\mathbf{A}\mathbf{f}\le\mathbf{b}$ & \cref{eq:cm-qp-std} \\
摩擦金字塔 & $\pm\mathbf{t}_{i,j}^\top\mathbf{f}_i\le\mu_i\mathbf{n}_i^\top\mathbf{f}_i$ & \cref{eq:cm-qp-friction} \\
无源性 & $\dot V\le\mathbf{u}^\top\mathbf{y}$ & \cref{eq:cm-passivity} \\
波变量 & $u=\tfrac{1}{\sqrt{2b}}(F+b\dot x),\ v=\tfrac{1}{\sqrt{2b}}(F-b\dot x)$ & \cref{eq:cm-wave} \\
\bottomrule
\end{tabular}
\end{table}

\subsection*{知识点总表}
\begin{table}[H]\centering\small
\caption{本章知识点与难度}
\label{tab:cm-knowledge}
\begin{tabular}{p{4.6cm}p{7.0cm}c}
\toprule
知识点 & 核心要点 & 难度 \\
\midrule
闭链约束 & 多机 $+$ 刚体 $=$ 闭合运动链，运动相互约束 & $\bigstar\bigstar\bigstar$ \\
抓取矩阵 $\mathbf{G}$ & $\mathbf{w}_o=\mathbf{G}\mathbf{f}$，列 $[\mathbf{I}_3;\skewm{\mathbf{r}}_i]$，纯几何不含动力学 & $\bigstar\bigstar\bigstar$ \\
速度对偶 & $\dot{\mathbf{X}}=\mathbf{G}^\top\mathbf{v}_o$，由功守恒与力映射对偶 & $\bigstar\bigstar\bigstar$ \\
零空间维度 & $\dim\ker\mathbf{G}=cN-6$，量化内力自由度 & $\bigstar\bigstar\bigstar$ \\
接触模型分类 & 点接触（$c{=}3$）/软指（$c{=}4$）/刚性抓持（$c{=}6$） & $\bigstar\bigstar$ \\
内力/外力分解 & $\mathbf{f}=\mathbf{G}^{+}\mathbf{w}+\mathbf{E}\boldsymbol{\lambda}$，正交分解 & $\bigstar\bigstar\bigstar$ \\
零空间投影性质 & 幂等 $+$ 对称，运动力 $\perp$ 内力 & $\bigstar\bigstar\bigstar$ \\
内力基 $\mathbf{E}$ & $\ker\mathbf{G}$ 标准正交基，$\boldsymbol{\lambda}$ 是内力模式强度 & $\bigstar\bigstar\bigstar$ \\
非挤压伪逆 & 最小范数 $\ne$ 非挤压，内力须显式设计 & $\bigstar\bigstar\bigstar$ \\
leader 阻抗律 & leader 用阻抗领航物体 & $\bigstar\bigstar\bigstar$ \\
follower 力放大 & 对齐物体运动方向，放大 leader 力 & $\bigstar\bigstar\bigstar$ \\
意图估计 & 物体运动方向即 leader 意图 & $\bigstar\bigstar\bigstar$ \\
去中心化同构控制 & 导航 $+$ 顺从 $+$ 内力三分量，本地律 & $\bigstar\bigstar\bigstar$ \\
隐式通信 & 物体 $=$ 零延迟无源广播信道 & $\bigstar\bigstar\bigstar$ \\
三架构权衡谱 & 全集中—主从—去中心化，信息 vs 规模 & $\bigstar\bigstar\bigstar$ \\
物体阻抗律 & 阻抗加在物体质心，$6$ 维统一柔顺 & $\bigstar\bigstar\bigstar$ \\
物体阻抗 $\to$ 接触力管线 & 阻抗 $\to$ 所需旋量 $\to$ 分解 $\to$ 关节力矩 & $\bigstar\bigstar\bigstar$ \\
内力阻抗 & 内力也柔顺，与运动阻抗正交解耦 & $\bigstar\bigstar\bigstar$ \\
负载分配 QP & 等式（旋量）$+$ 不等式（摩擦/单边/力限） & $\bigstar\bigstar\bigstar$ \\
摩擦金字塔 & 二阶锥线性化为多面锥供 QP & $\bigstar\bigstar\bigstar$ \\
三分配策略 & 等分/加权伪逆/QP 取舍 & $\bigstar\bigstar\bigstar$ \\
无源性定义与判稳 & $\dot V\le\mathbf{u}^\top\mathbf{y}$，无源 $\Rightarrow$ 稳定 & $\bigstar\bigstar\bigstar\bigstar$ \\
互联封闭定理 & 无源子系统互联仍无源 & $\bigstar\bigstar\bigstar\bigstar$ \\
能量罐 & 变阻抗产能须从罐取能 & $\bigstar\bigstar\bigstar\bigstar$ \\
波变量 & 延迟信道无源化 & $\bigstar\bigstar\bigstar\bigstar$ \\
前沿解读 & 绳索/被动臂/零通信学习 & $\bigstar\bigstar\bigstar\bigstar$ \\
\bottomrule
\end{tabular}
\end{table}

% ════════════════════════════════════════════════════════════════
\section*{延伸阅读}
\addcontentsline{toc}{section}{延伸阅读}
\begin{itemize}
  \item \textbf{奠基与经典}：Orin \& Oh\cite{orin1986force}（闭链力分配一般公式，\cref{sec:cm-qp} 的源头）\pz{年份应为 1981、ASME JDSMC 103(2)；bib 键 orin1986force 年份/卷期待修}；Walker, Freeman \& Marcus\cite{walker1988internal}（非挤压伪逆，内力经典分析）\pz{年份应为 1989}；Khatib\cite{khatib1988object} 与操作空间统一方法\cite{khatib1987operational}（增广物体模型与运动/力解耦，\cref{sec:cm-internal,sec:cm-impedance} 理论母体）；Hogan 阻抗控制\cite{hogan1985impedance}（阻抗律奠基，\cref{sec:cm-impedance} 形式来源）。
  \item \textbf{物体级阻抗与抓取控制}：Wimb\"ock, Ott, Albu-Sch\"affer \& Hirzinger\cite{wimbock2012comparison}（物体级阻抗统一框架，多指手通用，\cref{sec:cm-impedance} 现代形态）；Caccavale 等\cite{caccavale2008six}（双臂六自由度物体阻抗）\pz{此条 2008 IEEE/ASME T-Mech 正确；但"对称构型(symmetric formulation)"应归 Uchiyama \& Dauchez(1988/1993)，非 Caccavale 原创}。
  \item \textbf{分布式与去中心化}：Verginis \& Dimarogonas（及 Zelazo）\cite{verginis2019cooperative}（刚度理论视角，抓取矩阵与刚度矩阵 range-nullspace 关系、能量最优内力）\pz{发表于 IEEE T-CNS(2023)，非 RA-L}；Culbertson, Slotine \& Schwager\cite{culbertson2021decentralized}（去中心化自适应，免负载惯性）；隐式通信框架\cite{wang2018collaborative}\rebuilt{wang2018collaborative 键的 Frontiers 2018 题对应论文实为 Bechlioulis \& Kyriakopoulos，非 Wang \& Schwager；待重置作者或改引 Wang \& Schwager 真实工作(2014 DARS / 2016 IJRR)}；Carey \& Werfel\cite{carey2021collective}（隐式协调$+$自适应柔顺集体运输）\rebuilt{"Force-ANTS 力放大 N-robot 系统"标签张冠李戴——Force-ANTS 属 Wang \& Schwager 2016 IJRR；此条真题为 \emph{Collective Transport of Unconstrained Objects via Implicit Coordination and Adaptive Compliance}}。
  \item \textbf{无源性与遥操作}：Anderson \& Spong\cite{anderson1989bilateral}（散射理论解决延迟）；Niemeyer \& Slotine\cite{niemeyer1991stable}（波变量）；Hannaford \& Ryu\cite{hannaford2002time}（时域无源性/能量罐）；Nu\~no, Basa\~nez \& Ortega\cite{nuno2011passivity}（无源遥操作综述）。
  \item \textbf{前沿系统（研究级）}：多足绳牵搬运\cite{sombolestan2022collaborative}\rebuilt{键 sombolestan2022collaborative 张冠李戴：该 cable-towed-load 标题实为 Yang 等(UC Berkeley) RA-L 2022；Sombolestan \& Nguyen 的真实相关作为 IROS 2023 刚体自适应力控（另文），待修}；Turrisi 等\cite{turrisi2024pacc}（PACC 被动臂搬运，IROS 2024）；Pandit 等\cite{pandit2025decplm}（decPLM 零通信学习搬运，已接收 ICRA 2026）。
\end{itemize}

\section*{本章与后续章节的关系}
\addcontentsline{toc}{section}{本章与后续章节的关系}
\begin{table}[H]\centering\small
\begin{tabular}{p{3.4cm}p{6.4cm}p{3.8cm}}
\toprule
章节 & 关系 & 本章铺垫的知识点 \\
\midrule
\cref{ch:mr-consensus} 共识与分布式优化 & wrench 一致性正是 ADMM 的 $\mathbf{x}_i=\mathbf{z}$ 在“力”上的具体化；本章反向复用其分布式协调与一致性约束 & ADMM、一致性约束、共识 \\
异构多机协同 & 本章是同构搬运基础；异构搬运按机型分配不同子任务/力，复用抓取矩阵 $+$ 负载分配 QP & 抓取矩阵、负载分配、leader--follower \\
双臂/多臂协同操作 & 双臂是 $N{=}2$ 刚性抓持的特例，直接用本章抓取矩阵与内力分解 & 内力/外力分解、物体阻抗 \\
多足协同 Loco-Manipulation & 移动 $+$ 搬运的耦合；本章力控 $+$ 编队 MPC 的融合 & 物体阻抗、去中心化、QP 力分配 \\
MARL 系列（\cref{ch:plan-safemarl} 等） & decPLM 式学习搬运的 MARL 视角；本章理论是解读学习策略的物理框架 & 去中心化、隐式通信、内力 \\
\bottomrule
\end{tabular}
\end{table}

\noindent\textbf{与单臂力控章节的交叉引用}（横向关联，已在正文反复使用）：操作空间动力学 $\leftrightarrow$ \cref{sec:cm-grasp}（雅可比 $\mathbf{J}$ 与抓取矩阵 $\mathbf{G}$ 同构，$\boldsymbol{\tau}=\mathbf{J}^\top\mathbf{F}$ 与 $\dot{\mathbf{X}}=\mathbf{G}^\top\mathbf{v}_o$ 同源）；力控导论（阻抗/导纳二分法）$\leftrightarrow$ \cref{sec:cm-impedance}（单臂阻抗律推广到物体质心）；经典力控（摩擦锥）$\leftrightarrow$ \cref{sec:cm-qp}（摩擦锥/金字塔在负载分配 QP 复用）；变阻抗无源性与碰撞安全 $\leftrightarrow$ \cref{sec:cm-passivity}（单臂能量罐推广到多机耦合）；腿足 MPC/WBC 力控 $\leftrightarrow$ \cref{sec:cm-qp}（足端力分配 QP 与多机负载分配 QP 属同一 QP 族；详见 \cref{sec:cm-qp} 该处 \pz 对"完全一致"的界定）；双臂协调力控 $\leftrightarrow$ \cref{sec:cm-grasp,sec:cm-internal,sec:cm-impedance}（同一套抓取矩阵理论在 $N{=}2$ 刚性抓持的展开，本章是其多机推广）。

\section*{故障排查手册}
\addcontentsline{toc}{section}{故障排查手册}
\begin{enumerate}
  \item \textbf{症状}：物体被越夹越紧/接触力电流飙升。\textbf{原因}：内力失控（无内力调节，或各末端独立阻抗期望不一致）。\textbf{对策}：检查是否用了物体阻抗而非各末端独立阻抗；打印内力分量 $\mathbf{f}_{\text{int}}=\mathbf{E}\boldsymbol{\lambda}$ 看是否单调增；确认期望位姿只在物体级定义、各末端由刚体几何导出（\cref{sec:cm-internal,sec:cm-impedance}）。
  \item \textbf{症状}：物体从某接触点滑脱。\textbf{原因}：该点切向力超摩擦锥/内力（夹紧）不足。\textbf{对策}：检查 QP 是否含摩擦锥约束（非伪逆）；打印各点 $\|\mathbf{f}_t\|$ 与 $\mu f_n$ 比值；增大内力夹紧强度或法向预紧（\cref{sec:cm-qp,sec:cm-internal}）。
  \item \textbf{症状}：物体平移正常、一旋转就发散/翻转。\textbf{原因}：抓取矩阵下块 $\skewm{\mathbf{r}}_i$ 符号错/$\mathbf{r}_i$ 方向反。\textbf{对策}：单元测试 $\skewm{\mathbf{r}}\mathbf{v}=\mathbf{r}\times\mathbf{v}$；确认 $\mathbf{r}_i=$ 接触点 $-$ 质心；用纯平移测试通过、纯旋转测试定位（\cref{sec:cm-grasp}）。
  \item \textbf{症状}：力分配 QP 报 infeasible。\textbf{原因}：物理不可行（物体太重/摩擦太小/布局差）。\textbf{对策}：临时放松 $f^{\max},\mu$ 看是否变可行；检查接触点是否撑开（$\mathbf{r}_i$ 几何）；降低运动加速度（减小 $\mathbf{w}_{\text{ext}}$）；加机器人（\cref{sec:cm-qp}）。
  \item \textbf{症状}：follower 力方向剧烈抖动。\textbf{原因}：用未滤波的物体速度算方向，低速时 $0/0$ 病态。\textbf{对策}：对物体速度低通滤波；速度模低于阈值设死区/冻结方向；用观测器估计平滑速度（\cref{sec:cm-leader}）。
  \item \textbf{症状}：加机器人后系统低频振荡/发散。\textbf{原因}：无源性被破坏（新控制器阻尼不足/变阻抗产能/通信延迟）。\textbf{对策}：验证新控制器无源（写储能函数查 $\dot V\le\mathbf{u}^\top\mathbf{y}$）；变阻抗是否挂能量罐；有通信延迟是否用波变量；增大阻尼 $\mathbf{D}_d$（\cref{sec:cm-passivity}）。
  \item \textbf{症状}：去中心化搬运物体原地打转/乱走。\textbf{原因}：各 agent “同构控制律”坐标系未对齐。\textbf{对策}：确认导航方向在共同参考系定义；检查各接触力在世界系是否真同向；验证顺从分量未被关闭（\cref{sec:cm-decentral}）。
  \item \textbf{症状}：物体被抬起后持续下坠/稳态低于期望。\textbf{原因}：漏重力旋量补偿。\textbf{对策}：检查 $\mathbf{w}_{\text{ext}}$ 是否含 $-\mathbf{w}_{\text{grav}}$；悬停时验证接触力竖直分量之和 $\approx$ 物体重（\cref{sec:cm-impedance}）。
  \item \textbf{症状}：变阻抗策略一调刚度就失稳。\textbf{原因}：刚度增大注入能量，破坏无源性。\textbf{对策}：确认能量罐 $T_{\min}>0$ 且罐空时冻结产能；检查回收的耗散能量符号/来源正确；限幅 $T_{\max}$（\cref{sec:cm-passivity}）。
\end{enumerate}
